%%%%%%%%%%%%%%%%%%%%%%%%%%%%%%%%%%%%%%%%%%%%%%%%%%%%%%%%%%%%%%%%%%%%%%%%%%%%%%
%  Marketing Management (IM345AT) -- Volume 4 of 5
%  UNIT IV : Web Design and Mobile Marketing
%%%%%%%%%%%%%%%%%%%%%%%%%%%%%%%%%%%%%%%%%%%%%%%%%%%%%%%%%%%%%%%%%%%%%%%%%%%%%%

\documentclass{MMTextbook}

\unitnumber{IV}
\unittitle{Web Design and Mobile Marketing}
\unitsubtitle{The hub of the campaign ecosystem, and the device the customer
never puts down}
\unithours{07 Hours}
\volumeno{4}

\makeindex[intoc=true,columns=2,title=Index]

\hypersetup{
  pdftitle={Marketing Management IM345AT -- Unit IV: Web Design and Mobile
             Marketing},
  pdfsubject={Semester IV, Department of Industrial Engineering and Management,
              RV College of Engineering},
  pdfkeywords={web design, user-centred design, website optimization, mobile
               marketing, push notifications, IM345AT, marketing management}
}

\begin{document}

%%%%%%%%%%%%%%%%%%%%%%%%%%%%%%%%%%%%%%%%%%%%%%%%%%%%%%%%%%%%%%%%%%%%%%%%%%%%%%
\frontmatter
\pagestyle{mmfront}

\halftitlepage
\maintitlepage

\colophonpage{%
\textbf{\coursetitle\ ---\ The Study Text}\\[2pt]
Volume 4 of 5 \textbullet\ Unit IV: Web Design and Mobile Marketing\\[10pt]
%
Prepared for the course \coursetitle\ (\coursecode), Semester IV, of the
\coursedept, \courseinstitute.\\[10pt]
%
\textbf{Sources.} The syllabus, course outcomes, assessment rubrics and
reference-book list reproduced in the front matter are transcribed from the
official course document. The previous-year questions in Appendix~B are
transcribed from RVCE Continuous Internal Evaluation and Semester End
Examination papers for this course, and where an official scheme of valuation
was released the model answers follow it. The explanatory text is compiled from
classroom lecture notes and, where those were incomplete, from the prescribed
reference texts.\\[10pt]
%
\textbf{Typography.} Set in Nimbus Roman with URW Gothic display headings.
Composed in \LaTeX{} using the \texttt{MMTextbook} class written for this
series. The visual design is derived from \emph{The Legrand Orange Book}
template by Mathias Legrand, with modifications by Vel, distributed through
LaTeXTemplates.com under a Creative Commons BY-NC-SA 4.0 licence; the class
re-implements that design language on standard CTAN packages. All figures are
drawn in \TikZ{} and \texttt{pgfplots}.\\[10pt]
%
\textbf{Status.} Unofficial student-compiled study material. Not issued,
endorsed or verified by the Department or by RV College of Engineering.\\[10pt]
%
2026 edition.}

% ---------------------------------------------------------------- preface
\frontchapter{Preface to Volume 4}

Unit~IV covers the two places where every other channel finally lands. Volume~1
established that a digital campaign is a hub with spokes, and that the hub is a
quality company website. This unit is about building that hub properly --- and
about the device on which most visitors will actually reach it.

The first half is web design, and its governing idea is that design is not
decoration. The website is the conversion destination that every spoke feeds, so
a design decision is a revenue decision. The material accordingly runs from
principles to process: what web design is, what optimising a site means, how a
site is actually published, how to let the user rather than the design team drive
the decisions, how design and copy work as one system, and how to turn website
metrics into decisions rather than dashboards.

The second half is mobile marketing, and it opens with a distinction the
examiners return to more often than any other in this unit: mobile
\emph{advertising} is a paid subset of mobile \emph{marketing}. Getting that one
line right, and being able to expand it into a table, is worth more marks in
Unit~IV than any other single fact.

Chapters~1 to~5 cover web design; Chapters~6 to~8 cover mobile marketing. The
chapter sequence follows the syllabus phrase order exactly, and Appendix~A
records the mapping.

Two notes on the shape of this volume.

First, some material here deliberately repeats Volume~1. Krug's first law, the
seven usability attributes and the landing-page essentials appear in Unit~I under
``Optimize Website UX \& Landing Pages'' and again in Unit~IV under ``Design
Principles''. Both unit descriptors name them, so both volumes carry them; the
treatment here is the design-side treatment rather than the content-marketing
side.

Second, like Volume~3 this volume needs no appendix of off-syllabus material.
Every topic in the source notes maps to a phrase in the unit descriptor.

\vspace{12pt}
\noindent\emph{The whole of the course syllabus is reproduced overleaf. The
unit covered by this volume is Unit~IV.}

% ------------------------------------------------- conventions & syllabus
%%%%%%%%%%%%%%%%%%%%%%%%%%%%%%%%%%%%%%%%%%%%%%%%%%%%%%%%%%%%%%%%%%%%%%%%%%%%%%
%  mm-howtouse.tex
%  Editorial conventions for the series.  Shared, identical, across all five
%  volumes so that a reader who learns the apparatus once can use any volume.
%%%%%%%%%%%%%%%%%%%%%%%%%%%%%%%%%%%%%%%%%%%%%%%%%%%%%%%%%%%%%%%%%%%%%%%%%%%%%%

\frontchapter{How to Use This Study Text}

\section*{The shape of a volume}
\addcontentsline{toc}{section}{The shape of a volume}

Each of the five volumes in this series covers exactly one unit of the
IM345AT syllabus and is self-contained. Every volume opens with the official
syllabus for the whole course, so that the unit in hand can always be located
within it, and closes with a set of appendices carrying the concordance,
the question bank, the glossary and a quick-reference sheet.

The numbered chapters of each volume map one-to-one onto the topic phrases of
that unit's official syllabus descriptor. The order of the chapters is the
order of the syllabus. Appendix~A sets out that mapping explicitly, phrase by
phrase, so coverage can be audited rather than assumed.

\section*{The five panel types}
\addcontentsline{toc}{section}{The five panel types}

Material that is not continuous prose is set in one of five panels. Each has a
fixed meaning throughout the series.

\begin{definitionbox}[the term being defined]
A formal definition, worth reproducing close to verbatim in an answer script.
The term appears in the panel heading and again in the index.
\end{definitionbox}

\begin{keybox}[High Yield]
A result, list or distinction that carries disproportionate examination weight.
Where a topic has recurred across papers, the recurrence is recorded as a
bracketed tag. \pyq{CIE-I, 3 Apr 2025; CIE-I, 16 Apr 2026}
\end{keybox}

\begin{exambox}[Model Answer]
An answer written at the length and in the register the marking scheme
rewards. Where an official scheme of valuation was released, the panel follows
it.
\end{exambox}

\begin{examplebox}[Worked Example]
A concrete illustration --- a named company, a named persona, a worked
calculation --- rather than a general statement.
\end{examplebox}

\begin{warnbox}
A common and costly error: a distinction usually collapsed, a term usually
misused, or an answer that looks complete and is not.
\end{warnbox}

A sixth panel, in green, sets out a framework or process as a single
sequence:

\begin{frameworkbox}[Framework]
Step one \ra\ step two \ra\ step three \ra\ back to step one.
\end{frameworkbox}

\section*{How previous-year questions are recorded}
\addcontentsline{toc}{section}{How previous-year questions are recorded}

\begin{keybox}[The Bracket Convention]
A past-paper appearance is recorded \textbf{only} as a bracketed source tag
attached to the material, never as a sentence of narrative. So the text reads
\begin{quote}\sffamily\footnotesize
Push versus pull marketing. \pyq{CIE-I, 3 Apr 2025; CIE-I, 16 Apr 2026; SEE, Jun/Jul 2025 --- 2 M}
\end{quote}
and never ``this was asked in 2025 and again in 2026''. The tag carries the
paper, the date and, where the mark allocation is known, the marks. Multiple
appearances are separated by semicolons.
\end{keybox}

\noindent The abbreviations used inside tags are:

\vspace{4pt}
{\sffamily\small
\begin{tabular}{@{}P{0.14\textwidth}P{0.80\textwidth}@{}}
\toprule[1.1pt]
\rlab{CIE-I} & First Continuous Internal Evaluation test\\
\rlab{CIE-II} & Second Continuous Internal Evaluation test\\
\rlab{CIE-III} & Third Continuous Internal Evaluation test\\
\rlab{SEE} & Semester End Examination\\
\rlab{2 M} & Mark allocation of the question as printed on the paper\\
\bottomrule[1.1pt]
\end{tabular}}

\vspace{10pt}

Appendix~B of each volume reproduces every past-paper question traced to that
unit, in full, grouped by mark weight, each carrying the same bracketed tag.

\section*{Cross-references, figures and the index}
\addcontentsline{toc}{section}{Cross-references, figures and the index}

Cross-references are given as \S\,chapter.section --- so \S\,1.4 is the fourth
section of Chapter~1. Every figure in this series is drawn in \TeX{} rather
than imported as an image, so the diagrams print at full resolution at any
size and use the same typeface as the surrounding text. Figures and tables are
listed after the table of contents.

Terms set in \textbf{\textit{bold italic}} on first use are indexed. The index
at the back of each volume is a real back-of-book index built from those
terms, not a repeat of the contents list.

\section*{Status of this text}
\addcontentsline{toc}{section}{Status of this text}

\begin{warnbox}[Unofficial Document]
This is a student-compiled study text. It is not issued, endorsed or verified
by the Department of Industrial Engineering and Management or by RV College of
Engineering. The syllabus reproduced in the front matter, the assessment
rubrics and the previous-year questions are transcribed from official
documents; the explanatory prose around them is not official. Where the source
lecture notes were incomplete, gaps have been filled from the prescribed
reference texts listed in the front matter. Always cross-check against the
official syllabus and your own class notes.
\end{warnbox}

%%%%%%%%%%%%%%%%%%%%%%%%%%%%%%%%%%%%%%%%%%%%%%%%%%%%%%%%%%%%%%%%%%%%%%%%%%%%%%
%  mm-syllabus.tex
%  Verbatim reproduction of the official IM345AT course document.
%  Shared, byte-identical, across all five volumes of the series.
%%%%%%%%%%%%%%%%%%%%%%%%%%%%%%%%%%%%%%%%%%%%%%%%%%%%%%%%%%%%%%%%%%%%%%%%%%%%%%

\frontchapter{Official Syllabus: IM345AT}

\noindent
{\sffamily\small\color{slatelight}%
Reproduced from the official course document issued by the Department of
Industrial Engineering and Management, RV College of Engineering. The text of
the five unit descriptors below is the authority against which every chapter of
this series is written; nothing outside it appears in the numbered chapters of
any volume.\par}

\vspace{10pt}

\begin{center}
\sffamily\small
\begin{tabular}{@{}P{0.238\textwidth}P{0.238\textwidth}P{0.218\textwidth}P{0.218\textwidth}@{}}
\toprule[1.1pt]
\multicolumn{4}{@{}l}{\bfseries\normalsize Semester: IV \quad\textbullet\quad MARKETING MANAGEMENT}\\
\multicolumn{4}{@{}l}{\color{slatelight}Category: Professional Core Course (Theory)}\\
\midrule
\rlab{Course Code} & IM345AT & \rlab{CIE Marks} & 100 Marks\\
\rlab{Credits (L:T:P)} & 3:0:0 & \rlab{SEE Marks} & 100 Marks\\
\rlab{Total Hours} & 45 L & \rlab{SEE Duration} & 3.00 Hours\\
\bottomrule[1.1pt]
\end{tabular}
\end{center}

\vspace{6pt}

% ---------------------------------------------------------------- unit blocks
\newcommand{\syllunit}[3]{%
  \begin{tcolorbox}[mmbase, colback=white, colframe=ocre!45, boxrule=0.7pt,
      sharp corners, fontupper=\small,
      fonttitle=\sffamily\bfseries\small\color{white},
      coltitle=white, colbacktitle=#3,
      title={#1},
      attach boxed title to top left={xshift=-0.7pt, yshift=-0.7pt},
      boxed title style={sharp corners, boxrule=0pt, colback=#3},
      top=16pt]
    #2
  \end{tcolorbox}}

\syllunit{UNIT--I \hfill 07 Hrs}{%
\textbf{Introduction to Digital Marketing:} Principles of Digital Marketing;
Digital Marketing Channels; Tools to Create Buyer Persona; Competitor Research
Tools, Website Analysis Tools, etc.\par
\vspace{3pt}
\textbf{Content Marketing:} Content Marketing Concepts \& Strategies; Planning,
Creating, Distributing \& Promoting Content; Optimize Website UX \& Landing
Pages; Measure Impact; Metrics \& Performance; Using Content Research for
Opportunities, etc.}{ocre}

\syllunit{UNIT--II \hfill 08 Hrs}{%
\textbf{Social Media Marketing:} Introduction; Major Social Media Platforms for
Marketing; Developing Data-driven Audience \& Campaign Insights; Social Media
for Business; Creation \& Optimization of Social Media Campaigns, etc.\par
\vspace{3pt}
\textbf{Search Engine Optimization:} Search Engine Optimization Fundamentals;
Keywords and SEO Content Plan; SEO \& Business Objectives; Writing SEO Content;
On-site \& off-site SEO; Optimize Organic Search Ranking, etc.}{ocre}

\syllunit{UNIT--III \hfill 07 Hrs}{%
\textbf{Web Analytics \& Google Analytics:} Google Analytics Tools; Web
Analytics Tools, etc.\par
\vspace{3pt}
\textbf{E-mail Marketing:} Effective E-mail Campaigns; E-mail Plan; E-mail
Marketing Campaign Analysis; Measuring Conversions \& keeping up, etc.}{ocre}

\syllunit{UNIT--IV \hfill 07 Hrs}{%
\textbf{Web Design:} Web design, optimization of websites; Publishing a basic
website; User-centered Design and Website Optimization; Design Principles and
Website Copy; Website Metrics \& Developing Insight, etc.\par
\vspace{3pt}
\textbf{Mobile Marketing:} Difference between mobile advertising and marketing,
utilizing mobile marketing for sales promotions, online applications,
etc.}{ocre}

\syllunit{UNIT--V \hfill 07 Hrs}{%
\textbf{Conversion Optimization:} What is AIDAS and its role; website
optimization; what visitors want to see on the website; how to optimize key
element and increase the effect of landing on a particular page\par
\vspace{3pt}
\textbf{Digital Analytics:} Evolution of Digital Analytics, information about
end-to-end customer experience, analyst's influence on business, role as a
change agent, etc.}{ocre}

% ------------------------------------------------------------ course outcomes
\section*{Course Outcomes}
\addcontentsline{toc}{section}{Course Outcomes}

\noindent{\small After completing the course, the students will be able to:\par}
\vspace{4pt}

{\sffamily\small
\begin{longtable}{@{}P{0.07\textwidth}P{0.89\textwidth}@{}}
\toprule[1.1pt]
\rlab{CO1} & Differentiate the benefits drawn by updated marketing mix from
traditional marketing mix for effective marketing management there by to stay
competitive in today's global market-place.\\
\midrule
\rlab{CO2} & Develop an effective holistic marketing atmosphere to efficiently
face the challenges in dynamically changing market.\\
\midrule
\rlab{CO3} & Formulate a potential marketing plan to effectively reach the
targeted market segments, by delivering the value to targeted customers through
practicing sound marketing research.\\
\midrule
\rlab{CO4} & Create new channels to improvise marketing to achieve and maintain
competitive position in globalized market-place.\\
\bottomrule[1.1pt]
\end{longtable}}

% --------------------------------------------------------------- reference books
\section*{Prescribed Reference Books}
\addcontentsline{toc}{section}{Prescribed Reference Books}

{\small
\begin{enumerate}[leftmargin=2.1em]
\item \textit{Marketing Management}, Philip Kotler, Kevin Lane Keller, 15th
      Edition, 2016, Pearson, ISBN: 978-93-325-5718-5.
\item \textit{Digital Marketing --- Strategy, Implementation \& Practice}, Dave
      Chaffey, Fiona Ellis-Chadwick, 7th Edition, 2019, Pearson,
      ISBN: 9781292241623, 1292241624.
\item \textit{Marketing Research}, Donald S. Tull, Del I. Hawkins, 6th Edition,
      Prentice Hall India, ISBN: 8120309618.
\item \textit{Marketing Management --- A South Asian Perspective}, Philip
      Kotler, Kevin Lane Keller, Abraham Koshy, Mithileshwar Jha, 14th Edition,
      2013, Pearson, ISBN: 978-81-317-6716-0.
\item \textit{Marketing Research}, David A. Aaker, V. Kumar, George S. Day, 9th
      Edition, 2008, John Wiley \& Sons, ISBN: 978-265-1791-6.
\end{enumerate}}

% ------------------------------------------------------------------- rubrics
\frontchapter{Assessment Scheme}

\section*{Rubric for the Continuous Internal Evaluation (Theory)}
\addcontentsline{toc}{section}{Rubric for the Continuous Internal Evaluation}

{\sffamily\small
\begin{longtable}{@{}P{0.045\textwidth}P{0.80\textwidth}>{\raggedleft\arraybackslash}P{0.07\textwidth}@{}}
\toprule[1.1pt]
\rowcolor{ocre}\thd{\#} & \thd{Component} & \thd{Marks}\\
\midrule
\rlab{1.} & \textbf{Quizzes:} Quizzes will be conducted in online/offline mode.
TWO QUIZZES will be conducted \& each quiz will be evaluated for 10 marks. The
sum of two quizzes will be the final quiz marks. & \bfseries 20\\
\midrule
\rowcolor{tablealt}
\rlab{2.} & \textbf{Tests:} Students will be evaluated in test, descriptive
questions with different complexity levels (Revised Bloom's Taxonomy Levels:
Remembering, Understanding, Applying, Analyzing, Evaluating and Creating). TWO
tests will be conducted. Each test will be evaluated for 50 marks, adding up to
100 marks. Final test marks will be reduced to 40 marks. & \bfseries 40\\
\midrule
\rlab{3.} & \textbf{Experiential Learning:} Students will be evaluated for their
creativity and practical implementation of the problem. Case study-based
teaching learning (10), Program specific requirements (10), Video based
seminar/presentation/demonstration (20), adding up to 40. & \bfseries 40\\
\midrule
\multicolumn{2}{@{}l}{\bfseries MAXIMUM MARKS FOR THE CIE THEORY} & \bfseries 100\\
\bottomrule[1.1pt]
\end{longtable}}

\section*{Rubric for the Semester End Examination (Theory)}
\addcontentsline{toc}{section}{Rubric for the Semester End Examination}

{\sffamily\small
\begin{longtable}{@{}P{0.12\textwidth}P{0.72\textwidth}>{\raggedleft\arraybackslash}P{0.08\textwidth}@{}}
\toprule[1.1pt]
\rowcolor{ocre}\thd{Q. No.} & \thd{Contents} & \thd{Marks}\\
\midrule
\multicolumn{3}{@{}l}{\bfseries\color{slate} PART A}\\
\midrule
\rlab{1} & Objective type questions covering entire syllabus & \bfseries 20\\
\midrule
\multicolumn{3}{@{}l}{\bfseries\color{slate} PART B \quad\normalfont\itshape(Maximum of TWO sub-divisions only)}\\
\midrule
\rlab{2}      & Unit 1: (Compulsory)      & \bfseries 16\\
\rowcolor{tablealt}
\rlab{3 \& 4} & Unit 2: (Internal Choice) & \bfseries 16\\
\rlab{5 \& 6} & Unit 3: (Internal Choice) & \bfseries 16\\
\rowcolor{tablealt}
\rlab{7 \& 8} & Unit 4: (Internal Choice) & \bfseries 16\\
\rlab{9 \& 10}& Unit 5: (Internal Choice) & \bfseries 16\\
\midrule
\multicolumn{2}{@{}l}{\bfseries TOTAL} & \bfseries 100\\
\bottomrule[1.1pt]
\end{longtable}}

\vspace{8pt}

\begin{keybox}[What the SEE Structure Means for This Volume]
Part A draws twenty marks of objective questions from the \emph{entire}
syllabus, so no unit can be skipped. In Part B, Question~2 is compulsory and is
always set from Unit~I; Units~II to~V each carry a 16-mark internal-choice
pair. A candidate therefore answers Unit~I plus one question from each of the
remaining four units --- five full questions in all.
\end{keybox}


% ---------------------------------------------------------------- contents
\cleardoublepage
\pdfbookmark[0]{Contents}{toc}
\tableofcontents
\cleardoublepage
\listoffigures
\cleardoublepage
\listoftables

%%%%%%%%%%%%%%%%%%%%%%%%%%%%%%%%%%%%%%%%%%%%%%%%%%%%%%%%%%%%%%%%%%%%%%%%%%%%%%
\mainmatter
\pagestyle{mmmain}

%%%%%%%%%%%%%%%%%%%%%%%%%%%%%%%%%%%%%%%%%%%%%%%%%%%%%%%%%%%%%%%%%%%%%%%%%%%%%%
\chapter{Web Design and the Optimization of Websites}
%%%%%%%%%%%%%%%%%%%%%%%%%%%%%%%%%%%%%%%%%%%%%%%%%%%%%%%%%%%%%%%%%%%%%%%%%%%%%%

\syllabusstrip{Web Design: Web design, optimization of websites.}

\begin{learningoutcomes}
\begin{itemize}
\item What web design is, and why it is a revenue decision.
\item User interface and user experience --- scope, deliverables, marketing role
      and failure symptoms.
\item The basic pages a business website needs.
\item What website optimization is, why it matters, and its four types.
\item The eight optimisation techniques, with their concrete sub-actions.
\item Image optimization, and its four distinct contributions.
\item Testing and optimisation --- methods, what each finds, and the tools.
\end{itemize}
\end{learningoutcomes}

\section{What Web Design Is}

\begin{definitionbox}[Web Design]
\keytermas{Web design}{web design} involves creating visually appealing and
user-friendly websites. It includes layout design, colour selection, typography,
graphics, navigation and responsiveness.
\pyq{CIE-III, Jun 2026 --- 10 M}
\end{definitionbox}

\begin{keybox}[Why a Design Decision Is a Revenue Decision]
Web design is not decoration. In marketing terms the website is the \textbf{hub of
the entire digital campaign ecosystem} --- the round-the-clock portal and the
ultimate conversion destination that every other channel feeds. Every visitor
that SEO, paid search, social, e-mail and affiliates deliver arrives here. A
design flaw therefore taxes every channel simultaneously.
\end{keybox}

\section{User Interface and User Experience}

\begin{keybox}[High Yield]
``What is the role of user Interface and user Experience designs in web design
for marketing? How can businesses ensure a positive user Journey on their
websites?''
\pyq{SEE, Oct/Nov 2023, Q7a --- 8 M}
\end{keybox}

\begin{mmlongtable}{User interface against user experience}{tab:uiux}%
{@{}P{0.15\textwidth}P{0.39\textwidth}P{0.39\textwidth}@{}}
\rowcolor{ocre}\thd{} & \thd{User Interface (UI)} & \thd{User Experience (UX)}\\
\midrule
\endfirsthead
\rowcolor{ocre}\thd{} & \thd{User Interface (UI)} & \thd{User Experience (UX)}\\
\midrule
\endhead
\rlab{Definition} & The visual and interactive surface the user touches ---
layout, colours, typography, buttons, icons, spacing, states & The entire quality
of the user's journey --- how easy, efficient and satisfying it is to achieve
their goal\\
\rowcolor{tablealt}
\rlab{Scope} & What it looks like and how it responds & How it feels and whether
it works\\
\rlab{Deliverables} & Style guides, design systems, mockups, component libraries
& Personas, user journeys, wireframes, information architecture, usability test
findings\\
\rowcolor{tablealt}
\rlab{Marketing role} & Creates the first impression and signals brand
credibility within milliseconds & Determines whether the visitor converts,
returns and recommends\\
\rlab{Failure symptom} & The site looks dated or inconsistent; the brand appears
untrustworthy & Users cannot find what they need; high bounce, low conversion\\
\end{mmlongtable}

\begin{figure}[htbp]
\centering
\begin{tikzpicture}[font=\sffamily\footnotesize]
  % UX as the containing journey
  \draw[ocre,line width=1.1pt,rounded corners=6pt,fill=ocrepale]
    (0,0) rectangle (11.2,3.5);
  \node[anchor=north west,text=ocredark,font=\sffamily\bfseries\small]
    at (0.25,3.35) {USER EXPERIENCE --- the whole journey};

  % UI as one component within it
  \draw[deepblue,line width=1.1pt,rounded corners=5pt,fill=deepbluepale]
    (3.6,0.45) rectangle (7.6,2.35);
  \node[anchor=north,text=deepblue,font=\sffamily\bfseries\scriptsize,align=center]
    at (5.6,2.2) {USER INTERFACE\\ the surface};
  \node[anchor=north,text=slate,font=\sffamily\scriptsize,align=center,
        text width=3.6cm] at (5.6,1.5)
    {layout, colour, type, buttons, icons, spacing, states};

  % journey stages around it
  \node[anchor=center,text=slate,font=\sffamily\scriptsize,align=center,
        text width=2.9cm] at (1.75,1.4)
    {\textbf{before}\\ expectation set by the ad or search result};
  \node[anchor=center,text=slate,font=\sffamily\scriptsize,align=center,
        text width=2.9cm] at (9.45,1.4)
    {\textbf{after}\\ fulfilment, support, return, recommendation};

  \draw[-{Latex[length=2.4mm]},ocre,line width=1pt] (3.25,1.4) -- (3.55,1.4);
  \draw[-{Latex[length=2.4mm]},ocre,line width=1pt] (7.65,1.4) -- (7.95,1.4);

  \node[anchor=north,text=slatelight,font=\sffamily\scriptsize\itshape,
        align=center,text width=10.6cm] at (5.6,0.42)
    {a beautiful interface inside a broken journey still fails --- which is why
     UI is a component of UX and not a synonym for it};
\end{tikzpicture}
\caption[User interface within user experience]%
        {UI is one component of UX, not a synonym for it. The interface can be
         faultless while the journey --- expectation, fulfilment, support ---
         still loses the customer.}
\label{fig:uiux}
\end{figure}

\begin{exambox}[Ensuring a Positive User Journey --- Eight Actions]
Map the journey before designing; keep navigation shallow and conventional; make
the primary action obvious on every page; minimise the number of steps to the
goal; provide feedback for every action; design the error and empty states, not
just the happy path; make the site fast and responsive; and test with real users
rather than with colleagues.
\end{exambox}

\section{Website Structure --- The Basic Pages}

\begin{itemize}
\item \sfb{Home page} --- states what the business does, for whom, and what to do
      next, within five seconds.
\item \sfb{About page} --- credibility, team, history, credentials.
\item \sfb{Product and service pages} --- one page per offering, benefit-led,
      with proof and a call to action\ix{call to action}.
\item \sfb{Landing pages} --- single-purpose campaign destinations with
      navigation removed.
\item \sfb{Blog and resources} --- the SEO and content-marketing engine.
\item \sfb{Contact page} --- multiple contact routes, a map, hours, and a
      response-time expectation.
\item \sfb{Trust pages} --- testimonials, case studies, certifications, privacy
      policy, terms, refund policy.
\end{itemize}

\section{Website Optimization --- Definition and Importance}

\begin{definitionbox}[Website Optimization]
\keytermas{Website optimization}{website optimization} is the process of
improving website performance, speed, usability\ix{usability} and search visibility in order to
increase traffic and conversions.
\pyq{CIE-III, Jun 2026 --- 2 M}
\end{definitionbox}

\begin{exambox}[Importance --- Official Scheme, 10 Marks]
\pyq{CIE-III, Jun 2026 --- 10 M}
\begin{itemize}
\item Improves website loading speed.
\item Enhances user experience.
\item Increases search engine ranking.
\item Reduces bounce\ix{bounce} rate.
\item Improves conversion rate.
\item Supports mobile compatibility.
\end{itemize}

\smallskip
\textbf{Types of website optimization:} technical optimization; content
optimization; SEO optimization; mobile optimization.

\smallskip
\textbf{Benefits:} better customer engagement; increased sales and traffic;
improved brand image.
\end{exambox}

\section{The Optimization Techniques}

\begin{keybox}[The Recurring Scenario]
``A start-up company notices that visitors leave their website within a few
seconds. As a web designer, explain how you would optimize the website to improve
user engagement and conversion rate.''
\pyq{CIE-III, Jun 2026 --- 10 M}
\end{keybox}

\begin{enumerate}
\item \sfb{Improve website loading speed} --- compress images, minimise
      unnecessary scripts, use fast hosting, enable caching and a content
      delivery network, and minify CSS and JavaScript.
\item \sfb{Create a responsive design} --- ensure compatibility with mobile,
      tablet and desktop; design mobile-first.
\item \sfb{Improve navigation} --- use simple menus and a proper page structure;
      keep important pages within two or three clicks.
\item \sfb{Enhance content quality} --- attractive headings, meaningful visuals
      and clear information; remove needless words.
\item \sfb{Add strong call-to-action buttons} --- ``Buy Now'', ``Register'',
      ``Contact Us'' --- prominent, above the fold, and repeated.
\item \sfb{Improve SEO} --- proper keywords, unique title tags and meta
      descriptions, heading hierarchy, image ALT text.
\item \sfb{Reduce bounce rate} --- provide engaging, relevant content that matches
      the promise made by the referring advertisement or search result.
\item \sfb{Monitor website metrics} --- analyse visitor behaviour using analytics
      tools, and keep iterating.
\end{enumerate}

\begin{keybox}[How to Turn This Into a First-Class Answer]
Open with a \textbf{diagnosis step} --- heatmaps and session recordings to
establish \emph{why} visitors leave --- and close with \textbf{A/B testing} to
validate each fix. A list of eight techniques without a diagnosis at the front
and a validation at the back is a checklist; with them, it is a method.

\smallskip
Additional techniques for a fuller answer: A/B test headlines, calls to action
and layouts; simplify forms by removing non-essential fields; add trust signals
near the point of decision; implement exit-intent offers; use heatmaps and
session recordings to locate friction; and fix accessibility so the site works
for every visitor.
\end{keybox}

\section{Image Optimization}

\begin{exambox}[Model Answer --- 2 Marks]
``How does image optimization contribute to overall website performance?''
\pyq{Improvement CIE, 28 Aug 2024 --- 2 M}
\begin{itemize}
\item \sfb{Speed.} Reduces page weight and load time, which directly lowers
      bounce rate and improves Core Web Vitals\ix{Core Web Vitals}.
\item \sfb{Mobile.} Improves the experience on slower connections and limited
      data plans.
\item \sfb{SEO.} Descriptive file names and ALT text make images discoverable in
      image search and help the page rank.
\item \sfb{Accessibility.} ALT text is read by screen readers and displayed when
      images fail to load.
\end{itemize}
\end{exambox}

\noindent\textbf{Techniques:} compress files; use modern formats such as WebP and
AVIF; size images to their display dimensions; serve responsive images with
\texttt{srcset}; lazy-load below-the-fold images; and use a content delivery
network.

\section{Testing and Optimization}

\begin{keybox}[High Yield]
``Describe the importance of testing and optimization in web design for
Marketing. What tools and methods can be used?''
\pyq{SEE, Oct/Nov 2023, Q7b --- 8 M}
\end{keybox}

\begin{keybox}[Why Testing Matters]
Design opinions are unreliable predictors of user behaviour. Testing replaces
assumption with evidence, de-risks redesigns, produces compounding conversion
gains on all future traffic, and reveals device- and segment-specific failures
that are invisible in aggregate data.
\end{keybox}

\begin{mmlongtable}{Testing methods, what they find, and the tools}{tab:testing}%
{@{}P{0.22\textwidth}P{0.36\textwidth}P{0.35\textwidth}@{}}
\rowcolor{ocre}\thd{Method} & \thd{What it finds} & \thd{Tools}\\
\midrule
\endfirsthead
\rowcolor{ocre}\thd{Method} & \thd{What it finds} & \thd{Tools}\\
\midrule
\endhead
\rlab{A/B and multivariate testing} & Which version performs better on a defined
goal & VWO, Optimizely, AB Tasty\\
\rowcolor{tablealt}
\rlab{Usability testing} & Where real users get stuck, and why & UserTesting,
Maze, moderated sessions\\
\rlab{Heatmaps, scroll and click maps} & What attracts attention and what is
ignored & Hotjar, Microsoft Clarity, Crazy Egg\\
\rowcolor{tablealt}
\rlab{Session recordings} & Hesitation, rage clicks, dead clicks, form
abandonment & Hotjar, Clarity, FullStory\\
\rlab{Analytics funnel analysis} & The exact step at which users drop out &
Google Analytics 4\\
\rowcolor{tablealt}
\rlab{Performance testing} & Speed and Core Web Vitals & PageSpeed Insights,
Lighthouse, GTmetrix, WebPageTest\\
\rlab{Cross-browser and device testing} & Rendering failures on specific clients
& BrowserStack, LambdaTest\\
\rowcolor{tablealt}
\rlab{Accessibility auditing} & Barriers for users with disabilities & WAVE, axe,
Lighthouse\\
\rlab{Technical SEO crawling} & Broken links, duplicate meta, redirect chains &
Screaming Frog, SEMrush Site Audit\\
\end{mmlongtable}

\begin{summarybox}
Web design creates visually appealing, user-friendly sites through layout,
colour, typography, graphics, navigation and responsiveness --- and because the
site is the hub every channel feeds, a design decision is a revenue decision. UI
is the surface; UX is the whole journey, and UI is a component of UX rather than a
synonym. A business site needs home, about, product, landing, blog, contact and
trust pages. Website optimization improves performance, speed, usability and
search visibility to raise traffic and conversions, in four types --- technical,
content, SEO and mobile. The eight techniques run from load speed and
responsiveness through navigation, content, calls to action, SEO and bounce
reduction to metric monitoring; framed properly they open with diagnosis and close
with validation. Image optimization contributes speed, mobile experience, SEO and
accessibility. Testing replaces opinion with evidence through split tests,
usability sessions, heatmaps, recordings, funnel analysis, performance,
cross-device, accessibility and crawl\ix{crawl} audits.
\end{summarybox}

%%%%%%%%%%%%%%%%%%%%%%%%%%%%%%%%%%%%%%%%%%%%%%%%%%%%%%%%%%%%%%%%%%%%%%%%%%%%%%
\chapter{Publishing a Basic Website}
%%%%%%%%%%%%%%%%%%%%%%%%%%%%%%%%%%%%%%%%%%%%%%%%%%%%%%%%%%%%%%%%%%%%%%%%%%%%%%

\syllabusstrip{Publishing a basic website.}

\begin{learningoutcomes}
\begin{itemize}
\item The nine official steps from requirement analysis to maintenance.
\item What each step involves in practice, including the pre-launch technical
      setup that the nine-step list compresses.
\end{itemize}
\end{learningoutcomes}

\begin{keybox}[The Recurring Scenario]
``A college plans to publish its departmental website for admissions and student
activities. Explain the complete process involved in designing and publishing the
website.''
\pyq{CIE-III, Jun 2026 --- 10 M}
The official scheme lists nine steps. Name all nine in order, then expand.
\end{keybox}

\section{The Nine Steps}

\begin{figure}[htbp]
\centering
\begin{tikzpicture}[font=\sffamily\footnotesize,
    st/.style={draw=ocre,line width=0.85pt,fill=ocrepale,rounded corners=3pt,
               minimum height=1.35cm,text width=1.62cm,align=center,
               font=\sffamily\bfseries\scriptsize,text=slate,inner sep=3pt}]
  \node[st] (a1) at (0,0)     {1\\ requirement analysis};
  \node[st] (a2) at (1.90,0)  {2\\ website design};
  \node[st] (a3) at (3.80,0)  {3\\ content development};
  \node[st] (a4) at (5.70,0)  {4\\ domain registration};
  \node[st] (a5) at (7.60,0)  {5\\ web hosting};
  \foreach \a/\b in {a1/a2,a2/a3,a3/a4,a4/a5}{
    \draw[-{Latex[length=1.9mm]},ocre,line width=0.85pt] (\a)--(\b);}

  \node[st] (a6) at (1.90,-2.2)  {6\\ website development};
  \node[st] (a7) at (3.80,-2.2)  {7\\ testing};
  \node[st] (a8) at (5.70,-2.2)  {8\\ publishing};
  \node[st,draw=greenok,fill=greenpale] (a9) at (7.60,-2.2) {9\\ maintenance};
  \foreach \a/\b in {a6/a7,a7/a8,a8/a9}{
    \draw[-{Latex[length=1.9mm]},ocre,line width=0.85pt] (\a)--(\b);}

  \draw[-{Latex[length=1.9mm]},ocre,line width=0.85pt]
    (a5.south) -- ++(0,-0.55) -| (a6.north);

  \draw[-{Latex[length=2.1mm]},greenok,line width=0.85pt,dashed]
    (a9.east) -- ++(0.85,0) -- ++(0,2.2) -| (a3.north);
  \node[anchor=west,text=greenok,font=\sffamily\scriptsize\itshape,
        align=left,text width=2.4cm] at (8.65,-1.1)
    {maintenance returns to content};
\end{tikzpicture}
\caption[The nine-step website publishing process]%
        {The nine steps. Maintenance is not an epilogue: it loops back into
         content development, because a departmental site that stops publishing
         notices has failed even though it is technically live.}
\label{fig:ninesteps}
\end{figure}

\begin{mmlongtable}{The nine official steps}{tab:publishsteps}%
{@{}P{0.22\textwidth}P{0.71\textwidth}@{}}
\rowcolor{ocre}\thd{Step} & \thd{Activity}\\
\midrule
\endfirsthead
\rowcolor{ocre}\thd{Step} & \thd{Activity}\\
\midrule
\endhead
\rlab{1. Requirement analysis} & Identify website objectives and user needs\\
\rowcolor{tablealt}
\rlab{2. Website design} & Create the homepage, admission page, faculty page and
event sections; produce wireframes and visual design\\
\rlab{3. Content development} & Add text, images, notices and contact details\\
\rowcolor{tablealt}
\rlab{4. Domain registration} & Purchase a domain name\\
\rlab{5. Web hosting} & Select a hosting service to store the website files\\
\rowcolor{tablealt}
\rlab{6. Website development} & Build using HTML, CSS and JavaScript, or a
content management system\\
\rlab{7. Testing} & Check website speed, links, responsiveness and security\\
\rowcolor{tablealt}
\rlab{8. Publishing} & Upload the website files to the hosting server\\
\rlab{9. Maintenance} & Regularly update notices and student information\\
\end{mmlongtable}

\section{Expanding Each Step}

\begin{description}
\item[Requirement analysis] Define the audience --- prospective students, current
students, parents, recruiters --- the primary goal (admission enquiries), and the
success metric.
\item[Information architecture] Map the sitemap and navigation before any visual
design begins.
\item[Design] Wireframe first, then apply the design principles of Chapter~4,
mobile-first.
\item[Content development] Write scannable, benefit-led copy; obtain and optimise
images; prepare downloadable documents.
\item[Domain registration] Choose a short, memorable, relevant name; register with
an accredited registrar; configure DNS.
\item[Hosting] Select shared, VPS, cloud or managed hosting based on expected
traffic; ensure an SSL certificate for HTTPS.
\item[Development] Build with semantic HTML for accessibility and SEO; integrate a
content management system so that non-technical staff can update notices.
\item[Pre-launch technical setup] Install analytics and Search Console, submit an
XML sitemap, verify that \texttt{robots.txt} does not block the site, and set up
backups.
\item[Testing] Functional testing of every link and form; cross-browser and
cross-device testing; speed testing; security testing; an accessibility audit;
and content proofreading.
\item[Publishing] Deploy files via SFTP or a version-controlled pipeline, point
DNS to the host, and verify HTTPS and redirects.
\item[Post-launch maintenance] Regular content updates, software and plugin
updates, security patching, backup verification, broken-link checks, and a monthly
analytics review.
\end{description}

\begin{warnbox}[The Step the Nine-Step List Hides]
The official list moves straight from development to testing, but there is a
\emph{pre-launch technical setup} between them that is where sites are most often
lost. Launching with analytics uninstrumented means the launch baseline is gone
for ever; launching with a staging \texttt{noindex} directive still in place means
the site is invisible to search engines. Name this step explicitly and the answer
reads as practitioner knowledge rather than recitation.
\end{warnbox}

\begin{summarybox}
Publishing a website runs nine steps: requirement analysis, website design,
content development, domain registration, web hosting, website development,
testing, publishing and maintenance. In practice, information architecture
precedes design, a pre-launch technical setup --- analytics, Search Console,
sitemap, \texttt{robots.txt} check, backups --- sits between development and
testing, and maintenance loops back into content rather than ending the process.
\end{summarybox}

%%%%%%%%%%%%%%%%%%%%%%%%%%%%%%%%%%%%%%%%%%%%%%%%%%%%%%%%%%%%%%%%%%%%%%%%%%%%%%
\chapter{User-Centred Design and Website Optimization}
%%%%%%%%%%%%%%%%%%%%%%%%%%%%%%%%%%%%%%%%%%%%%%%%%%%%%%%%%%%%%%%%%%%%%%%%%%%%%%

\syllabusstrip{User-centered Design and Website Optimization.}

\begin{learningoutcomes}
\begin{itemize}
\item What user-centred design is, and the two principles to name.
\item The seven principles of user-centred design.
\item The five-stage UCD process.
\item The ten questions website visitors arrive with, and how to answer each.
\end{itemize}
\end{learningoutcomes}

\section{What User-Centred Design Is}

\begin{definitionbox}[User-Centred Design]
\keytermas{User-centred design}{user-centred design} (UCD) is a design philosophy
and process in which the needs, goals, context and limitations of the end user
drive every decision at every stage, and in which the design is validated by real
users rather than by the design team's judgement.
\end{definitionbox}

\begin{exambox}[Two Key Principles --- 2 Marks]
\pyq{Improvement CIE, 4 Jun 2025 --- 2 M}
(1) Design based on an explicit understanding of the users, their tasks and their
environment. (2) Iterative design that is continuously refined through user
involvement and user-centred evaluation.
\end{exambox}

\section{The Seven Principles}

\begin{enumerate}
\item \sfb{Design is based on an explicit understanding of users, tasks and
      environments} --- research first, design second.
\item \sfb{Users are involved throughout design and development} --- not
      consulted once at the end.
\item \sfb{The design is driven and refined by user-centred evaluation} ---
      usability testing at every iteration.
\item \sfb{The process is iterative} --- design, test, learn, redesign.
\item \sfb{The design addresses the whole user experience}, not just the
      interface --- including content, support and post-purchase.
\item \sfb{The design team includes multidisciplinary skills} and perspectives.
\item \sfb{Accessibility and inclusivity} --- the design must work for users with
      differing abilities, devices and connection speeds.
\end{enumerate}

\section{The UCD Process}

\begin{figure}[htbp]
\centering
\begin{tikzpicture}[font=\sffamily\footnotesize]
  \def\R{3.05}
  \foreach \i/\lbl/\ang in {%
      1/{1. Understand\\ the context\\ of use}/108,
      2/{2. Specify user\\ requirements}/36,
      3/{3. Produce design\\ solutions}/-36,
      4/{4. Evaluate\\ against\\ requirements}/-108,
      5/{5. Iterate}/180}{
    \node[draw=ocre,line width=0.9pt,fill=ocrepale,rounded corners=4pt,
          align=center,text width=2.25cm,inner sep=5pt,
          font=\sffamily\bfseries\scriptsize,text=slate] (n\i) at (\ang:\R) {\lbl};}
  \foreach \a/\b in {1/2,2/3,3/4}{
    \draw[-{Latex[length=2.3mm]},ocre,line width=1pt] (n\a) to[bend left=14] (n\b);}
  \draw[-{Latex[length=2.3mm]},ocre,line width=1pt] (n4) to[bend left=14] (n5);
  \draw[-{Latex[length=2.3mm]},ocre,line width=1pt] (n5) to[bend left=14] (n1);
  \node[circle,draw=greenok,line width=1pt,fill=greenpale,
        align=center,text width=1.85cm,font=\sffamily\bfseries\scriptsize,
        text=greenok,inner sep=2pt] at (0,0) {validated by real users};
\end{tikzpicture}
\caption[The user-centred design process]%
        {The UCD cycle. Iteration continues after launch: the loop has no exit,
         only release points.}
\label{fig:ucd}
\end{figure}

\begin{enumerate}
\item \sfb{Understand the context of use} --- who the users are, what they are
      trying to do, on what device, under what constraints.
\item \sfb{Specify user requirements} --- translate research into concrete
      functional and content requirements.
\item \sfb{Produce design solutions} --- information architecture, wireframes,
      prototypes, then visual design.
\item \sfb{Evaluate against requirements} --- usability testing, heuristic
      evaluation, A/B testing\ix{A/B testing}, analytics review.
\item \sfb{Iterate} until the design meets the user requirements --- then continue
      iterating after launch.
\end{enumerate}

\section{What Website Visitors Actually Look For}

\begin{keybox}[High Yield]
``Discuss what website visitors typically look for when they land on a business
webpage. Suggest effective design and content strategies to meet those
expectations.''
\pyq{Improvement CIE, 4 Jun 2025 --- 10 M}
The strongest structure is to state each visitor \emph{question} and pair it with
the delivery mechanism, as below.
\end{keybox}

\begin{mmlongtable}{The ten questions visitors arrive with}{tab:visitorwants}%
{@{}P{0.32\textwidth}P{0.61\textwidth}@{}}
\rowcolor{ocre}\thd{What visitors want to know} & \thd{How to deliver it}\\
\midrule
\endfirsthead
\rowcolor{ocre}\thd{What visitors want to know} & \thd{How to deliver it}\\
\midrule
\endhead
\rlab{``Am I in the right place?''} & A clear headline stating what the business
does and for whom, visible without scrolling\\
\rowcolor{tablealt}
\rlab{``What's in it for me?''} & Benefit-led copy, not feature lists or company
history\\
\rlab{``Can I trust you?''} & Testimonials, ratings, client logos, certifications,
guarantees, real photographs, a physical address\\
\rowcolor{tablealt}
\rlab{``How much does it cost?''} & Transparent pricing, or at least a clear
pricing range; hidden pricing is a leading cause of abandonment\\
\rlab{``What do I do next?''} & One prominent, unambiguous call to action\\
\rowcolor{tablealt}
\rlab{``Can I find what I came for?''} & Simple navigation and a working search\\
\rlab{``Is it fast?''} & A page that loads in under about three seconds\\
\rowcolor{tablealt}
\rlab{``Does it work on my phone?''} & A fully responsive, mobile-first layout\\
\rlab{``How do I reach a human?''} & Visible contact options --- phone, chat,
e-mail, hours\\
\rowcolor{tablealt}
\rlab{``Is it safe?''} & HTTPS, recognised payment logos, a clear privacy and
returns policy\\
\end{mmlongtable}

\begin{exambox}[Two Elements on a Well-Optimised Landing Page --- 2 Marks]
\pyq{Improvement CIE, 4 Jun 2025 --- 2 M}
(1) A clear, benefit-led headline that matches the advertisement or link that
brought the visitor there, and (2) a single prominent call-to-action button. Also
acceptable: trust signals or testimonials; a hero image or product visual; a short
form.
\end{exambox}

\begin{summarybox}
User-centred design lets the user's needs, goals, context and limitations drive
every decision, and validates the result with real users rather than with the
design team. Its seven principles require explicit user understanding, continuous
user involvement, evaluation-driven refinement, iteration, whole-experience scope,
multidisciplinary teams, and accessibility. The process runs understand the
context, specify requirements, produce solutions, evaluate, iterate --- and does
not stop at launch. Visitors arrive with ten questions, from ``am I in the right
place'' to ``is it safe'', and a well-designed page answers each one at the point
it is asked.
\end{summarybox}

%%%%%%%%%%%%%%%%%%%%%%%%%%%%%%%%%%%%%%%%%%%%%%%%%%%%%%%%%%%%%%%%%%%%%%%%%%%%%%
\chapter{Design Principles and Website Copy}
%%%%%%%%%%%%%%%%%%%%%%%%%%%%%%%%%%%%%%%%%%%%%%%%%%%%%%%%%%%%%%%%%%%%%%%%%%%%%%

\syllabusstrip{Design Principles and Website Copy.}

\begin{learningoutcomes}
\begin{itemize}
\item The twelve core design principles, and the two to name.
\item Krug's first law of usability and the seven usability attributes.
\item How people actually use the web, and why they do not read.
\item Billboard design --- designing for scanning.
\item The principles of effective website copy.
\item How design principles and copy work as one system.
\end{itemize}
\end{learningoutcomes}

\section{The Core Design Principles}

\begin{mmlongtable}{The core principles of web design}{tab:designprinciples}%
{@{}P{0.22\textwidth}P{0.71\textwidth}@{}}
\rowcolor{ocre}\thd{Principle} & \thd{What it requires}\\
\midrule
\endfirsthead
\rowcolor{ocre}\thd{Principle} & \thd{What it requires}\\
\midrule
\endhead
\rlab{Simplicity} & Remove everything that does not serve the user's goal. Every
added element competes with the call to action\\
\rowcolor{tablealt}
\rlab{Consistency} & The same colours, fonts, button styles, spacing and
terminology on every page, so that users do not relearn the interface\\
\rlab{Visual hierarchy} & Make important elements visually prominent in
proportion to their importance --- size, weight, colour and position\\
\rowcolor{tablealt}
\rlab{Balance and alignment} & Ordered use of a grid; alignment creates the
perception of quality\\
\rlab{Contrast} & Sufficient contrast between text and background for readability
and accessibility\\
\rowcolor{tablealt}
\rlab{Whitespace} & Deliberate empty space that groups related items and gives
the eye rest\\
\rlab{Typography} & Readable body size of 16 pixels or more, limited font
families, generous line height, restrained line length\\
\rowcolor{tablealt}
\rlab{Colour} & A limited, purposeful palette; one accent colour reserved for
actions\\
\rlab{Responsiveness} & The layout adapts to mobile, tablet and desktop ---
mobile-first by default\\
\rowcolor{tablealt}
\rlab{Accessibility} & Alt text, keyboard navigation, colour-independent meaning,
adequate tap targets\\
\rlab{Navigation clarity} & Conventional, shallow, labelled in the user's
language and not in internal jargon\\
\rowcolor{tablealt}
\rlab{Speed} & Performance is a design constraint, not an engineering
afterthought\\
\end{mmlongtable}

\begin{exambox}[Two Principles --- 2 Marks]
``Mention any two principles of website design.''
\pyq{CIE-III, Jun 2026 --- 2 M}
The scheme answer is \textbf{simplicity} and \textbf{consistency} --- these
principles improve readability and user experience.
\end{exambox}

\begin{exambox}[Two Aspects of Branding and Consistency --- 2 Marks]
\pyq{SEE, Oct/Nov 2023 --- 2 M}
(1) A consistent \textbf{visual identity} --- the same logo, colour palette,
typography, imagery style and button design applied on every page. (2) A
consistent \textbf{tone of voice and terminology} in all copy, including
microcopy, error messages and CTA labels, so that the brand reads as one coherent
personality.
\end{exambox}

\section{Krug's First Law, and the Seven Usability Attributes}

\begin{keybox}[Don't Make Me Think]
A person of average --- or below average --- ability must be able to work out how
to accomplish a task. The interface must be self-evident and self-explanatory.
The good news is that usability is not rocket science.
\end{keybox}

\begin{mmlongtable}{The seven usability attributes}{tab:usability4}%
{@{}P{0.18\textwidth}P{0.75\textwidth}@{}}
\rowcolor{ocre}\thd{Attribute} & \thd{Test question}\\
\midrule
\endfirsthead
\rowcolor{ocre}\thd{Attribute} & \thd{Test question}\\
\midrule
\endhead
\rlab{Useful}     & Does it do what people actually need?\\
\rowcolor{tablealt}
\rlab{Learnable}  & Can they figure it out?\\
\rlab{Memorable}  & Do they have to relearn it on every visit?\\
\rowcolor{tablealt}
\rlab{Effective}  & Does it get the job done?\\
\rlab{Efficient}  & Does it do so in a reasonable amount of time?\\
\rowcolor{tablealt}
\rlab{Desirable}  & Do people want it?\\
\rlab{Delightful} & Is it enjoyable to use?\\
\end{mmlongtable}

\section{How People Actually Use the Web}

There is a gap between how designers imagine people use websites --- reading
carefully, like a novel --- and how they actually do, which the heatmap\ix{heatmap} reveals:
\keyterm{scanning}, \keyterm{satisficing} and muddling through. Satisficing means
choosing the first option that plausibly works rather than searching for the best
one.

\noindent Why don't people read instructions? Three reasons.

\begin{enumerate}
\item It isn't important to them --- they just want it to work.
\item They are on a mission, and know they don't need to read everything.
\item If they find something that works, they stick to it rather than looking for
      a better way.
\end{enumerate}

\section{Billboard Design --- Design for Scanning, Not Reading}

\begin{itemize}
\item \sfb{Follow visual hierarchy} --- make important elements visually
      prominent.
\item \sfb{Omit needless words} --- the art of not writing for the web. Happy talk
      must die; unnecessary instructions must die.
\item \sfb{Make clickables obvious} --- it must be instantly clear what can be
      clicked, to prevent frustration. Competition is one click away.
\item \sfb{Provide breadcrumbs and street signs} --- design clear, intuitive
      navigation so that users always know where they are and how to get back.
\end{itemize}

\section{Principles of Effective Website Copy}

\begin{itemize}
\item \sfb{Lead with the benefit, not the feature.} Users buy outcomes, not
      specifications.
\item \sfb{Write for scanning} --- short paragraphs, meaningful subheadings,
      bullet lists, bold key phrases.
\item \sfb{Front-load the value} --- put the conclusion in the first line, not the
      last.
\item \sfb{Use the customer's language}, taken from reviews, support tickets and
      search queries --- not internal jargon.
\item \sfb{Be specific} --- ``ships in 24 hours'' beats ``fast delivery''.
\item \sfb{Use active voice and the second person} --- ``you get'', not
      ``customers are provided with''.
\item \sfb{Omit needless words.} Happy talk --- welcome messages that say nothing
      --- must go.
\item \sfb{One idea per section}, with a heading that states that idea.
\item \sfb{Make CTA copy explicit and active} --- ``Get my free trial''
      outperforms ``Submit''.
\item \sfb{Answer objections in the copy where they arise}, rather than hiding
      them in an FAQ.
\end{itemize}

\section{Design Principles and Copy as One System}

\begin{keybox}[High Yield]
``Discuss the relationship between design principles and website copy. How can
visual design elements and well-crafted content work together to enhance user
experience and achieve website objectives?''
\pyq{Improvement CIE, 28 Aug 2024 --- 10 M}
\end{keybox}

Design and copy are not sequential layers; they are a single system. Copy without
design is unreadable, and design without copy is empty.

\begin{mmlongtable}{How each design principle serves the copy}{tab:designcopy}%
{@{}P{0.20\textwidth}P{0.73\textwidth}@{}}
\rowcolor{ocre}\thd{Design principle} & \thd{How it serves the copy}\\
\midrule
\endfirsthead
\rowcolor{ocre}\thd{Design principle} & \thd{How it serves the copy}\\
\midrule
\endhead
\rlab{Visual hierarchy} & Determines the order in which the copy is read; the
headline must be visually dominant because it carries the core message\\
\rowcolor{tablealt}
\rlab{Typography} & Controls whether the copy is read at all --- line length,
size and spacing govern readability\\
\rlab{Whitespace} & Groups related copy and gives the reader cognitive rest,
increasing comprehension\\
\rowcolor{tablealt}
\rlab{Contrast} & Ensures the copy is legible, and directs attention to the CTA
text\\
\rlab{Consistency} & Means the same tone of voice and the same button language
appear everywhere, building trust\\
\rowcolor{tablealt}
\rlab{Imagery} & Demonstrates visually what the copy claims verbally --- proof,
not decoration\\
\rlab{Responsiveness} & Ensures the copy retains its hierarchy on a small screen,
where most users read it\\
\end{mmlongtable}

\begin{examplebox}[Effective Integration]
A \sfb{hero section} where a benefit headline, a supporting sub-line, a product
image and one accent-coloured CTA button occupy the top of the page with nothing
competing.

A \sfb{pricing table} where visual weight highlights the recommended plan while
the copy names the buyer it suits.

A \sfb{testimonial block} where a customer photograph and a short pull-quote in
large type do together what neither could do alone.
\end{examplebox}

\begin{warnbox}[The Two Characteristic Failures]
A beautiful layout with vague copy that never says what the company sells --- or
strong copy set in 11-pixel grey text on a white background, which nobody reads.

\smallskip
Each failure is a whole discipline working without the other.
\end{warnbox}

\begin{summarybox}
Twelve principles govern web design --- simplicity, consistency, visual
hierarchy, balance and alignment, contrast, whitespace, typography, colour,
responsiveness, accessibility, navigation clarity and speed --- of which
simplicity and consistency are the two to name. Krug's first law is don't make me
think, and usability is tested against seven attributes. People scan, satisfice
and muddle through rather than reading, so pages must be designed as billboards.
Effective copy leads with benefits, is written for scanning, front-loads value,
uses the customer's language, is specific and active, omits needless words and
answers objections where they arise. Design and copy are one system: hierarchy
sets reading order, typography decides whether copy is read, whitespace aids
comprehension, contrast directs attention, consistency builds trust, imagery
supplies proof and responsiveness preserves hierarchy on the screen most people
use.
\end{summarybox}

%%%%%%%%%%%%%%%%%%%%%%%%%%%%%%%%%%%%%%%%%%%%%%%%%%%%%%%%%%%%%%%%%%%%%%%%%%%%%%
\chapter{Website Metrics and Developing Insight}
%%%%%%%%%%%%%%%%%%%%%%%%%%%%%%%%%%%%%%%%%%%%%%%%%%%%%%%%%%%%%%%%%%%%%%%%%%%%%%

\syllabusstrip{Website Metrics \& Developing Insight.}

\begin{learningoutcomes}
\begin{itemize}
\item The key website metrics, and the four to name if only four are asked for.
\item How metrics help optimise for marketing effectiveness.
\item The progression from data to decision --- what makes a metric an insight.
\item User behaviour analysis and its impact on optimisation.
\end{itemize}
\end{learningoutcomes}

\section{Key Website Metrics}

\begin{mmlongtable}{Key website metrics}{tab:webmetrics}%
{@{}P{0.20\textwidth}P{0.34\textwidth}P{0.39\textwidth}@{}}
\rowcolor{ocre}\thd{Metric} & \thd{What it measures} & \thd{What it tells you}\\
\midrule
\endfirsthead
\rowcolor{ocre}\thd{Metric} & \thd{What it measures} & \thd{What it tells you}\\
\midrule
\endhead
\rlab{Traffic (users, sessions)} & Volume of visitors & Whether acquisition is
working\\
\rowcolor{tablealt}
\rlab{Bounce rate} & Percentage leaving after one page without interaction &
Relevance and first-impression quality\\
\rlab{Average session duration} & Time spent per visit & Depth of engagement\\
\rowcolor{tablealt}
\rlab{Pages per session} & Pages viewed per visit & Whether users explore or
leave\\
\rlab{Conversion rate} & Percentage completing the desired action & Whether the
site produces business value\\
\rowcolor{tablealt}
\rlab{Page load time and Core Web Vitals} & Speed and stability & Technical health
and its effect on abandonment\\
\rlab{Traffic sources} & Channel mix & Which marketing investments work\\
\rowcolor{tablealt}
\rlab{New against returning visitors} & Retention split & Whether the site earns
repeat visits\\
\rlab{Exit pages} & Where sessions end & Which page is losing people\\
\rowcolor{tablealt}
\rlab{Cart abandonment rate} & Percentage starting but not completing checkout &
Friction in the purchase flow\\
\rlab{Cost per acquisition} & Spend divided by conversions & Efficiency of the
whole funnel\\
\rowcolor{tablealt}
\rlab{Goal value and revenue} & Money produced & The bottom line\\
\end{mmlongtable}

\begin{exambox}[The Four-Metric Answer --- 2 Marks]
``List any four key metrics used to measure the success of a website.''
\pyq{Improvement CIE, 28 Aug 2024 --- 2 M}
Name \textbf{traffic, bounce rate, conversion rate and average session
duration}. These four cover acquisition, relevance, outcome and engagement
respectively, which reads as a considered selection rather than an arbitrary
list.
\end{exambox}

\section{How Metrics Help Optimise for Marketing Effectiveness}

\begin{exambox}[Model Answer --- 2 Marks]
\pyq{Improvement CIE, 4 Jun 2025 --- 2 M}
\begin{itemize}
\item They \sfb{locate the problem precisely} --- a high bounce rate on one
      landing page is actionable, whereas ``the site isn't working'' is not.
\item They \sfb{rank the fixes by impact} --- the step losing the most users is
      fixed first.
\item They \sfb{prove or disprove changes} through before-and-after comparison and
      A/B testing.
\item They \sfb{reallocate budget} towards the channels producing conversions
      rather than traffic.
\item They \sfb{reveal segment-specific failures} --- for example a mobile-only
      checkout bug hidden by desktop averages.
\item They \sfb{set realistic targets} by grounding goals in historical
      baselines.
\end{itemize}
\end{exambox}

\section{Developing Insight --- From Data to Decision}

\begin{keybox}[A Metric Is Not an Insight]
The progression runs: \textbf{data \ra{} metric \ra{} insight \ra{} decision
\ra{} action \ra{} measured result}. A dashboard stops at the second step, which
is why dashboards on their own change nothing.
\end{keybox}

\begin{figure}[htbp]
\centering
\begin{tikzpicture}[font=\sffamily\footnotesize,
    stp/.style={draw=ocre,line width=0.85pt,fill=ocrepale,rounded corners=3pt,
                minimum height=1.05cm,text width=2.05cm,align=center,
                font=\sffamily\scriptsize,text=slate,inner sep=3pt}]
  \node[stp] (s1) at (0,0)      {\textbf{1. Observe}\\ checkout completion 22\%
     against a 45\% benchmark};
  \node[stp] (s2) at (2.35,0)   {\textbf{2. Segment}\\ the gap is entirely on
     mobile; desktop is 47\%};
  \node[stp] (s3) at (4.70,0)   {\textbf{3. Diagnose}\\ recordings show
     abandonment at the address form};
  \node[stp] (s4) at (7.05,0)   {\textbf{4. Hypothesise}\\ nine fields is too
     many for a mobile keyboard};
  \node[stp] (s5) at (9.40,0)   {\textbf{5. Test}\\ A/B test the shortened
     form};
  \node[stp,draw=greenok,fill=greenpale] (s6) at (11.75,0)
     {\textbf{6. Decide}\\ implement the winner and re-measure};
  \foreach \a/\b in {s1/s2,s2/s3,s3/s4,s4/s5,s5/s6}{
    \draw[-{Latex[length=1.9mm]},ocre,line width=0.85pt] (\a)--(\b);}
  \draw[-{Latex[length=2.1mm]},greenok,line width=0.85pt,dashed]
    (s6.south) -- ++(0,-0.7) -| (s1.south);
  \node[anchor=north,text=greenok,font=\sffamily\scriptsize\itshape]
    at (5.85,-0.75) {the measured result becomes the next observation};
\end{tikzpicture}
\caption[From data to decision]%
        {The six moves that turn a metric into an insight. Steps~2 and~3 are the
         ones usually skipped, and they are the ones that make the hypothesis
         worth testing.}
\label{fig:datatodecision}
\end{figure}

\begin{enumerate}
\item \sfb{Observe} --- the checkout completion rate is 22\%, against an industry
      benchmark of 45\%.
\item \sfb{Segment} --- the gap is entirely on mobile; desktop is at 47\%.
\item \sfb{Diagnose qualitatively} --- session recordings show mobile users
      abandoning at the address form.
\item \sfb{Hypothesise} --- ``the nine-field address form is too long for a mobile
      keyboard; auto-fill and postcode lookup will reduce abandonment.''
\item \sfb{Test} --- run an A/B test on the new form.
\item \sfb{Decide and implement} the winner.
\item \sfb{Re-measure} and feed the result back into the next hypothesis.
\end{enumerate}

\section{User Behaviour Analysis}

\begin{exambox}[Model Answer --- 2 Marks]
``What is the impact of user behavior analysis on website optimization?''
\pyq{SEE, Oct/Nov 2023 --- 2 M}
User behaviour analysis --- heatmaps, scroll maps, click maps, session recordings
and funnel reports --- shows what users actually do rather than what they say or
what the business assumes.

\smallskip
\textbf{Its impact:} it identifies friction points precisely; it reveals ignored
content and invisible calls to action; it exposes rage clicks and dead clicks on
non-functional elements; it shows how far down the page users actually scroll,
which determines what must be placed above the fold; and it generates the
hypotheses that A/B testing then validates. Without it, optimisation is
guesswork.
\end{exambox}

\begin{summarybox}
Website success is measured on traffic, bounce rate, session duration, pages per
session, conversion rate, load time and Core Web Vitals, traffic sources, new
against returning visitors, exit pages, cart abandonment, cost per acquisition and
revenue --- and if only four are wanted, traffic, bounce rate, conversion rate and
session duration cover acquisition, relevance, outcome and engagement. Metrics
help by locating problems precisely, ranking fixes by impact, proving changes,
reallocating budget, exposing segment-specific failures and grounding targets. A
metric becomes an insight only through observation, segmentation, qualitative
diagnosis, hypothesis, test and decision. User behaviour analysis supplies the
diagnosis that makes hypotheses worth testing.
\end{summarybox}

%%%%%%%%%%%%%%%%%%%%%%%%%%%%%%%%%%%%%%%%%%%%%%%%%%%%%%%%%%%%%%%%%%%%%%%%%%%%%%
%  Unit IV, second half: Mobile Marketing (Chapters 6-8)
%%%%%%%%%%%%%%%%%%%%%%%%%%%%%%%%%%%%%%%%%%%%%%%%%%%%%%%%%%%%%%%%%%%%%%%%%%%%%%

\chapter{Mobile Advertising and Mobile Marketing}

\syllabusstrip{Mobile Marketing: Difference between mobile advertising and
marketing.}

\begin{learningoutcomes}
\begin{itemize}
\item What mobile marketing is, and the scale of mobile media consumption.
\item The difference between mobile advertising and mobile marketing --- the
      single most repeated question in this unit.
\item Mobile web marketing against mobile app marketing.
\item The benefits of mobile marketing, and its five real limitations.
\item The seven characteristics of a successful mobile campaign.
\item Geo-targeting and geo-fencing.
\end{itemize}
\end{learningoutcomes}

\section{What Mobile Marketing Is}

\begin{definitionbox}[Mobile Marketing]
\keytermas{Mobile marketing}{mobile marketing} is a multi-channel digital
marketing strategy aimed at reaching a target audience on their smartphones,
tablets and other mobile devices --- through mobile websites, e-mail, SMS and
MMS, social media, apps and push notifications.
\end{definitionbox}

Mobile is disrupting the way people engage with brands. Everything that can be
done on a desktop is now available on a mobile device --- from opening an e-mail
to visiting a website to reading content, all through a small screen. About
\textbf{80\%} of internet users own a smartphone, and mobile platforms host up to
\textbf{60\%} of digital media consumption.

\section{Mobile Advertising Against Mobile Marketing}

\begin{keybox}[The Most Repeated Question in Unit IV]
The same distinction, set in three different phrasings.
\pyq{CIE-III, Jun 2026; Improvement CIE, 4 Jun 2025; Improvement CIE, 28 Aug 2024 --- 2 M each}
\end{keybox}

\begin{exambox}[The One-Line Answer]
Mobile \textbf{advertising} focuses mainly on paid promotional messages delivered
to mobile devices. Mobile \textbf{marketing} includes overall customer engagement
through SMS, apps, social media and mobile websites.

\smallskip
In one line: \textbf{mobile advertising is a paid subset of the broader mobile
marketing strategy.}
\end{exambox}

\begin{figure}[htbp]
\centering
\begin{tikzpicture}[font=\sffamily\footnotesize]
  % outer set: mobile marketing
  \draw[ocre,line width=1.1pt,rounded corners=7pt,fill=ocrepale]
    (0,0) rectangle (11.2,4.35);
  \node[anchor=north west,text=ocredark,font=\sffamily\bfseries\small]
    at (0.3,4.18) {MOBILE MARKETING --- the whole strategy};
  \node[anchor=north west,text=slate,font=\sffamily\scriptsize,align=left,
        text width=4.4cm] at (0.3,3.55)
    {apps \ \textbullet\ \ mobile websites\\ SMS and MMS \ \textbullet\ \ push
     notifications\\ mobile social \ \textbullet\ \ mobile e-mail\\ QR codes \
     \textbullet\ \ location-based\\ mobile commerce};
  \node[anchor=north west,text=slatelight,font=\sffamily\scriptsize\itshape,
        align=left,text width=4.4cm] at (0.3,1.5)
    {paid \emph{and} organic\\ continuous, two-way\\ measured across the whole
     funnel};

  % inner subset: mobile advertising
  \draw[deepblue,line width=1.1pt,rounded corners=6pt,fill=deepbluepale]
    (5.55,0.55) rectangle (10.7,3.8);
  \node[anchor=north,text=deepblue,font=\sffamily\bfseries\scriptsize]
    at (8.12,3.65) {MOBILE ADVERTISING};
  \node[anchor=north,text=slatelight,font=\sffamily\scriptsize\itshape]
    at (8.12,3.28) {a paid subset --- one tactic};
  \node[anchor=north,text=slate,font=\sffamily\scriptsize,align=center,
        text width=4.6cm] at (8.12,2.88)
    {banner and display ads \ \textbullet\ \ in-app ads\\ mobile search ads \
     \textbullet\ \ interstitials\\ video ads \ \textbullet\ \ rewarded video};
  \node[anchor=north,text=slatelight,font=\sffamily\scriptsize\itshape,
        align=center,text width=4.6cm] at (8.12,1.55)
    {always paid \ \textbullet\ \ campaign-bound\\ largely one-way\\ measured on
     impressions, clicks, CTR};
\end{tikzpicture}
\caption[Mobile advertising as a subset of mobile marketing]%
        {The containment relation is the answer. Mobile advertising sits
         \emph{inside} mobile marketing: it is one paid tactic within a strategy
         that also runs organic, continuous and two-way activity.}
\label{fig:mobsubset}
\end{figure}

\begin{mmlongtable}{Mobile advertising against mobile marketing}{tab:advsmark}%
{@{}P{0.15\textwidth}P{0.38\textwidth}P{0.40\textwidth}@{}}
\rowcolor{ocre}\thd{Basis} & \thd{Mobile Advertising} & \thd{Mobile Marketing}\\
\midrule
\endfirsthead
\rowcolor{ocre}\thd{Basis} & \thd{Mobile Advertising} & \thd{Mobile Marketing}\\
\midrule
\endhead
\rlab{Scope} & Narrow --- one tactic & Broad --- the whole strategy\\
\rowcolor{tablealt}
\rlab{Nature} & Always paid & Paid and organic\\
\rlab{Channels} & Banner advertisements, in-app advertisements, interstitials,
mobile search advertisements, video advertisements & Apps, mobile websites, SMS
and MMS, push notifications, social, e-mail, QR codes, location-based marketing,
mobile commerce\\
\rowcolor{tablealt}
\rlab{Objective} & Immediate visibility, clicks and conversions & Long-term
engagement, relationship and retention\\
\rlab{Duration} & Campaign-bound & Continuous\\
\rowcolor{tablealt}
\rlab{Interaction} & Largely one-way & Two-way and interactive\\
\rlab{Measurement} & Impressions, clicks, click-through rate, cost per install &
Full-funnel: engagement, retention, lifetime value\\
\end{mmlongtable}

\section{Mobile Web Marketing Against Mobile App Marketing}

\begin{keybox}[High Yield]
``List any two differences between Mobile Web marketing and Mobile App
marketing.''
\pyq{SEE, Oct/Nov 2023 --- 2 M}
The two strongest pairs are \emph{access and reach}, and \emph{push and device
access}.
\end{keybox}

\begin{mmlongtable}{Mobile web against mobile app marketing}{tab:webvsapp}%
{@{}P{0.18\textwidth}P{0.37\textwidth}P{0.38\textwidth}@{}}
\rowcolor{ocre}\thd{Basis} & \thd{Mobile Web Marketing} & \thd{Mobile App Marketing}\\
\midrule
\endfirsthead
\rowcolor{ocre}\thd{Basis} & \thd{Mobile Web Marketing} & \thd{Mobile App Marketing}\\
\midrule
\endhead
\rlab{Access} & Through a browser; no installation required & Requires download
and installation\\
\rowcolor{tablealt}
\rlab{Reach} & Broader --- anyone with a browser and a link & Narrower --- only
those who install\\
\rlab{Discoverability} & Found through search engines (SEO) & Found through app
store optimisation (ASO)\\
\rowcolor{tablealt}
\rlab{Engagement depth} & Lower; session-based & Much higher; repeat usage and
habit formation\\
\rlab{Push notifications} & Limited, via web push & Full push notification
capability\\
\rowcolor{tablealt}
\rlab{Offline use} & Generally unavailable & Possible\\
\rlab{Device features} & Limited access to camera, GPS, contacts & Full access to
device hardware and sensors\\
\rowcolor{tablealt}
\rlab{Cost} & Lower to build and maintain & Higher --- separate platform builds\\
\end{mmlongtable}

\section{Benefits of Mobile Marketing}

\begin{itemize}
\item \sfb{Reach} --- mobile phones reach people about 15\% more than internet
      marketing does, and mobile advertising draws more clicks.
\item \sfb{Accessibility} --- the mobile device is always at hand; mobile
      advertising follows people everywhere.
\item \sfb{Time factor} --- people are available on mobile around the clock, which
      is not possible with desktops.
\item \sfb{Cost} --- mobile advertising costs much less than other types of
      advertising, allowing more advertising for the same spend.
\item \sfb{Personalised} --- the message can be tailored; people find mobile
      messages more intimate than other internet marketing methods.
\item \sfb{Reaches a broader market} --- smartphones and tablets are cheaper,
      smaller and more portable than PCs, so people previously excluded by
      financial, geographical or technological barriers are now online.
\item \sfb{Instantaneous results} --- we always carry our phones, and they are
      usually switched on, so the message is received the moment it is sent.
\item \sfb{Easy to work with} --- creating elements for mobile is simpler and less
      costly than for desktop; promotions such as downloadable coupons can be
      stored and redeemed without an internet connection.
\item \sfb{Convenient to use} --- the small screen limits content, which forces
      creators to keep it basic and simple; simpler content also adapts better
      across mobile platforms.
\item \sfb{Tracking response} --- user response can be tracked almost
      instantaneously, and phone numbers work well as unique identifiers because
      people keep them longer than e-mail addresses, which helps build accurate
      buyer personas.
\item \sfb{Huge viral potential} --- mobile content is easily shared, producing a
      domino effect and extra exposure at no additional cost.
\item \sfb{Mass communication made easy} --- more people own phones than
      computers, and mobile allows geo-targeting through GPS and Bluetooth as well
      as demographic targeting.
\item \sfb{Micro-blogging benefits} --- platforms such as Instagram and X place
      influence in the hands of ordinary people, so anyone can act as an
      influencer.
\item \sfb{Mobile payment} --- secure mobile payment environments let users
      purchase and pay bills without physical currency.
\end{itemize}

\section{Limitations}

\begin{warnbox}[The Five Real Constraints]
\begin{itemize}
\item \sfb{Platforms too diverse} --- mobile devices have no single standard;
      screen sizes vary and platforms use different operating systems and
      browsers, so one campaign for all of them is difficult.
\item \sfb{Privacy issues} --- users want their privacy respected, so marketers
      must offer clear opt-out instructions.
\item \sfb{Navigation on a mobile phone} --- small screens and no mouse make
      navigation difficult even with touch, so many advertisements go untouched
      because reviewing each one is tedious.
\item \sfb{Ad fatigue and intrusiveness} --- push notifications and interstitials
      are easily overused, driving uninstalls.
\item \sfb{Regulatory constraints} --- SMS and push marketing are governed by
      consent regulations, including do-not-disturb registries in India and GDPR
      in the European Union.
\end{itemize}
\end{warnbox}

\section{The Seven Characteristics of a Successful Mobile Campaign}

Derived from the Mobile Marketing Association's interviews with members and its
awards analysis.

\begin{enumerate}
\item \sfb{Think ``mobile first''.} The most successful campaigns are developed
      with a mobile-first mentality; starting from a mobile perspective gives the
      campaign a stronger foundation. Consider mobile the \emph{connective tissue}
      to all media, since it supports and strengthens every channel.
\item \sfb{Leverage multi-screen usage patterns.} Smartphone and tablet users work
      across screens --- watching a match on television while using the app for
      other scores, or posting while watching a show. Brands leverage cross-screen
      usage to capture attention.
\item \sfb{Utilise the full spectrum of mobile tools and applications.} Mobile
      isn't just a mobile website \emph{or} a display campaign \emph{or} an app
      --- it is a mobile website \emph{and} a display campaign \emph{and} an app,
      and much more.
\item \sfb{Integrate the mobile campaign into the traditional campaign.} One
      manufacturer generated 39 million total views by integrating its game-time
      app with a concurrent televised championship campaign. Mobile campaigns
      should never be produced in a silo.
\item \sfb{Create a campaign that works across multiple screens.} Tablet users are
      often at the \emph{top} of the sales funnel\ix{sales funnel} doing in-depth initial research,
      while smartphone users are often at the \emph{bottom}, ready to purchase.
\item \sfb{Leverage every phase of the sales funnel.} Mobile visitors use it for
      search and discovery as well as for purchase and brand connection.
\item \sfb{Test your way to success.} Mobile is digital by nature, making it
      perfect for tracking and measuring. The second half of the battle is sifting
      through those insights and adjusting to improve future performance.
\end{enumerate}

\begin{keybox}[The Closing Idea]
Mobile isn't simply a new sales channel or marketing tool. It is a revolutionary
new medium that is transforming not only the way consumers connect with brands,
but how they make purchases and stay engaged with them.
\end{keybox}

\section{Geo-Targeting}

\begin{definitionbox}[Geo-Targeting]
\keytermas{Geo-targeting}{geo-targeting} in mobile marketing is the delivery of
content, offers or advertisements to users based on their geographic location,
determined through GPS, Wi-Fi, Bluetooth beacons, cell-tower triangulation or IP
address.
\pyq{SEE, Jun/Jul 2025 --- 2 M}
\end{definitionbox}

\begin{exambox}[The Role of Geo-Targeting]
It makes messages locally relevant and timely --- sending an offer when the user
is physically near the store (\keytermas{geo-fencing}{geo-fencing}); enabling
``near me'' discovery and click-to-map directions; supporting location-specific
pricing, language and inventory; driving footfall to physical outlets; and
improving ROI by eliminating spend on users outside the serviceable area.
\end{exambox}

\begin{summarybox}
Mobile marketing is a multi-channel strategy reaching users on mobile devices
through websites, e-mail, SMS, social, apps and push notifications; about 80\% of
internet users own a smartphone and mobile carries up to 60\% of digital media
consumption. Mobile advertising is the paid subset of that strategy --- narrow,
always paid, campaign-bound, one-way and measured on clicks --- while mobile
marketing is broad, paid and organic, continuous, two-way and measured across the
whole funnel. Mobile web reaches more people with less depth; apps reach fewer
people with far more depth, full push capability, device access and offline use.
The benefits run from reach, accessibility and round-the-clock availability
through cost, personalisation and viral potential to mobile payment; the
limitations are platform fragmentation, privacy, small-screen navigation, ad
fatigue and consent regulation. Seven characteristics mark a successful campaign,
beginning with mobile-first thinking and ending with continuous testing.
Geo-targeting delivers location-relevant offers and geo-fencing triggers them at a
boundary.
\end{summarybox}

%%%%%%%%%%%%%%%%%%%%%%%%%%%%%%%%%%%%%%%%%%%%%%%%%%%%%%%%%%%%%%%%%%%%%%%%%%%%%%
\chapter{Utilizing Mobile Marketing for Sales Promotions}
%%%%%%%%%%%%%%%%%%%%%%%%%%%%%%%%%%%%%%%%%%%%%%%%%%%%%%%%%%%%%%%%%%%%%%%%%%%%%%

\syllabusstrip{Utilizing mobile marketing for sales promotions.}

\begin{learningoutcomes}
\begin{itemize}
\item Push notifications --- what they are, how they serve sales promotion, and
      the eight strategies that make them work.
\item Ten mobile sales-promotion techniques.
\item The mobile advertising formats, including click-to-call and click-to-map.
\end{itemize}
\end{learningoutcomes}

\section{Push Notifications}

\begin{definitionbox}[Push Notification]
A \keytermas{push notification}{push notification} is a short message sent by an
app directly to a user's device screen, appearing even when the app is not open.
\end{definitionbox}

\begin{keybox}[High Yield]
``Discuss briefly on Push Notification in Mobile Marketing. Explain how
businesses can use them to engage and retain Mobile App users? Provide examples
of effective push notification strategies.''
\pyq{SEE, Oct/Nov 2023, Q8a --- 8 M; Improvement CIE, 28 Aug 2024 --- 2 M}
\end{keybox}

\subsection{How push notifications enhance sales promotions}

\begin{itemize}
\item \sfb{Immediacy} --- a flash sale can be announced to the entire installed
      base within seconds.
\item \sfb{Urgency} --- countdown and limited-stock messaging drives immediate
      action.
\item \sfb{Personalisation} --- notifications triggered by browsing or purchase
      history: ``the jacket you viewed is now 30\% off''.
\item \sfb{Cart recovery} --- abandoned-cart reminders sent minutes after
      abandonment.
\item \sfb{Geo-triggered offers} --- a coupon delivered when the user enters a
      defined radius around a store.
\item \sfb{Re-engagement} --- win-back messages to users who have not opened the
      app in 30 days.
\item \sfb{Loyalty and replenishment} --- points-expiry alerts and reorder
      reminders for consumables.
\end{itemize}

\subsection{Eight effective push notification strategies}

\begin{enumerate}
\item \sfb{Segment before sending} --- never broadcast to the entire base;
      irrelevant pushes drive uninstalls.
\item \sfb{Time it correctly} --- send within the user's active hours in their own
      time zone.
\item \sfb{Cap frequency} --- set an explicit limit per user per week.
\item \sfb{Write with a single clear benefit and a verb}, within roughly 40
      characters for the title.
\item \sfb{Use rich media} --- images and action buttons materially raise
      engagement.
\item \sfb{Deep-link the notification} straight to the relevant screen, never to
      the home screen.
\item \sfb{Ask permission at the right moment} --- request notification consent
      after the user has experienced value, not on first launch.
\item \sfb{A/B test} copy, timing and creative, and measure opt-out rate as
      carefully as click rate.
\end{enumerate}

\begin{warnbox}[The Metric Most Often Ignored]
Push campaigns are usually judged on open rate alone. The \textbf{opt-out rate} is
the more important number, because an opt-out is permanent: a campaign that lifts
opens by 5\% while pushing 2\% of the base to disable notifications has destroyed
future reach to buy present clicks.
\end{warnbox}

\section{Mobile Sales-Promotion Techniques}

\begin{mmlongtable}{Ten mobile sales-promotion techniques}{tab:mobpromo}%
{@{}P{0.24\textwidth}P{0.69\textwidth}@{}}
\rowcolor{ocre}\thd{Technique} & \thd{How it works}\\
\midrule
\endfirsthead
\rowcolor{ocre}\thd{Technique} & \thd{How it works}\\
\midrule
\endhead
\rlab{SMS campaigns} & High open rate, often above 90\%; best for time-critical
offers and verified coupon codes\\
\rowcolor{tablealt}
\rlab{Mobile coupons} & Downloadable or in-app codes redeemable in store or
online; can be used offline\\
\rlab{In-app exclusive offers} & Discounts available only through the app, driving
installs and app engagement\\
\rowcolor{tablealt}
\rlab{Flash sales via push} & Time-boxed offers announced by push to create
urgency\\
\rlab{Location-based offers} & Geo-fenced coupons triggered near a store\\
\rowcolor{tablealt}
\rlab{QR codes} & Bridge print, packaging and outdoor media to a mobile landing
page or offer\\
\rlab{Gamification} & Spin-the-wheel, scratch cards and loyalty streaks to drive
repeat opens\\
\rowcolor{tablealt}
\rlab{Mobile wallet passes} & Offers saved to a device wallet with
location-triggered reminders\\
\rlab{Referral programmes} & One-tap share of a referral code from within the
app\\
\rowcolor{tablealt}
\rlab{Mobile-exclusive early access} & App users get the sale hours before the
website\\
\end{mmlongtable}

\section{Mobile Advertising Formats}

\begin{figure}[htbp]
\centering
\begin{tikzpicture}[font=\sffamily\footnotesize,
    ph/.style={draw=slate!60,line width=1pt,rounded corners=4pt,
               minimum width=2.35cm,minimum height=3.6cm,fill=white},
    fmtlabel/.style={anchor=north,align=center,text width=2.5cm,
                font=\sffamily\bfseries\scriptsize,text=ocredark},
    sub/.style={anchor=north,align=center,text width=2.6cm,
                font=\sffamily\scriptsize,text=slate}]

  % --- click to call
  \node[ph] (p1) at (0,0) {};
  \fill[ocrepale] (-1.05,0.95) rectangle (1.05,1.65);
  \draw[ocre,line width=0.8pt] (-1.05,0.95) rectangle (1.05,1.65);
  \node[font=\sffamily\scriptsize,text=ocredark] at (0,1.3) {Call now};
  \fill[rulegrey] (-1.05,-0.2) rectangle (1.05,0.55);
  \fill[rulegrey] (-1.05,-1.4) rectangle (1.05,-0.6);
  \node[fmtlabel] at (0,-2.05) {CLICK-TO-CALL};
  \node[sub] at (0,-2.5) {a tappable number; lifts CTR by 6--8\%};

  % --- interstitial
  \node[ph] (p2) at (2.95,0) {};
  \fill[deepbluepale] (1.9,-1.55) rectangle (4.0,1.75);
  \draw[deepblue,line width=0.8pt] (1.9,-1.55) rectangle (4.0,1.75);
  \node[font=\sffamily\scriptsize,text=deepblue,align=center] at (2.95,0.4)
    {full-screen\\ ad};
  \node[font=\sffamily\scriptsize,text=slatelight] at (3.78,1.6) {\texttimes};
  \node[fmtlabel] at (2.95,-2.05) {INTERSTITIAL};
  \node[sub] at (2.95,-2.5) {appears within an app; click through or close};

  % --- click to map
  \node[ph] (p3) at (5.90,0) {};
  \fill[greenpale] (4.85,-0.35) rectangle (6.95,1.75);
  \draw[greenok,line width=0.8pt] (4.85,-0.35) rectangle (6.95,1.75);
  \draw[greenok,line width=0.7pt] (5.0,0.5) -- (6.2,1.2) -- (6.8,0.3);
  \fill[warnred] (6.2,1.2) circle (0.09);
  \node[font=\sffamily\scriptsize,text=slate] at (5.90,-0.75) {nearest store};
  \node[fmtlabel] at (5.90,-2.05) {CLICK-TO-MAP};
  \node[sub] at (5.90,-2.5) {geo-targeted; opens a map, then contact details};

  % --- expandable
  \node[ph] (p4) at (8.85,0) {};
  \fill[ocre!25] (7.8,-1.55) rectangle (9.9,1.75);
  \draw[ocre,line width=0.8pt] (7.8,-1.55) rectangle (9.9,1.75);
  \draw[-{Latex[length=1.7mm]},ocredark,line width=0.7pt] (8.3,0.6) -- (7.95,0.95);
  \draw[-{Latex[length=1.7mm]},ocredark,line width=0.7pt] (9.4,0.6) -- (9.75,0.95);
  \draw[-{Latex[length=1.7mm]},ocredark,line width=0.7pt] (8.3,-0.4) -- (7.95,-0.75);
  \draw[-{Latex[length=1.7mm]},ocredark,line width=0.7pt] (9.4,-0.4) -- (9.75,-0.75);
  \node[font=\sffamily\scriptsize,text=slate,align=center] at (8.85,0.1)
    {rich media\\ expands};
  \node[fmtlabel] at (8.85,-2.05) {EXPANDABLE};
  \node[sub] at (8.85,-2.5) {canvas ad expands to cover the whole screen};
\end{tikzpicture}
\caption[Four mobile advertising formats]%
        {Four of the principal formats. Each trades intrusiveness against intent:
         click-to-call and click-to-map serve a user already looking for a
         business, while interstitials and expandables interrupt one who is not.}
\label{fig:mobformats}
\end{figure}

\begin{mmlongtable}{Mobile advertising formats}{tab:mobformats}%
{@{}P{0.24\textwidth}P{0.69\textwidth}@{}}
\rowcolor{ocre}\thd{Format} & \thd{Description}\\
\midrule
\endfirsthead
\rowcolor{ocre}\thd{Format} & \thd{Description}\\
\midrule
\endhead
\rlab{Click-to-Call} & Displays a phone number the user can tap to be connected
directly to a call centre or sales centre. Research shows click-to-call
advertisements drive a 6--8\% average increase in click-through rates\\
\rowcolor{tablealt}
\rlab{Interstitials} & Interactive advertisements that appear within an app; once
the user opens the app the interstitial displays, and the user can click through
to a landing page or close it\\
\rlab{Click-to-Map} & Geo-targeting sends messages to users in a specific
location; clicking the advertisement opens a map identifying the nearest store,
and clicking the map displays the contact information on the smartphone\\
\rowcolor{tablealt}
\rlab{Canvas and expandable advertisements} & When clicked, the advertisement
expands to cover the entire phone screen; it can be animated or incorporate rich
media to enhance the user experience\\
\rlab{Mobile search advertisements} & Paid results on mobile search pages, often
with location extensions\\
\rowcolor{tablealt}
\rlab{In-app banner and native advertisements} & Display units embedded in the
app's content flow\\
\rlab{Rewarded video} & The user watches a video in exchange for in-app currency
or content --- high completion rates\\
\end{mmlongtable}

\begin{keybox}[Do Not Send Mobile Advertising to a Static Page]
Mobile display advertisements are effective for building brand awareness\ix{brand awareness} and
generating clicks, leads and conversions. They need not link to a static landing
page --- linking to a \emph{dynamic} page or experience produces improved
engagement and higher conversions.
\end{keybox}

\begin{summarybox}
A push notification is a short message an app sends to the device screen even when
closed, and it serves sales promotion through immediacy, urgency,
personalisation, cart recovery, geo-triggered offers, re-engagement and loyalty
prompts. Eight strategies make it work: segment, time correctly, cap frequency,
write one benefit with a verb, use rich media, deep-link, ask permission after
value, and A/B test while watching opt-out as closely as opens. Ten promotion
techniques span SMS, coupons, in-app exclusives, push flash sales, geo-fenced
offers, QR codes, gamification, wallet passes, referrals and mobile-exclusive
early access. The advertising formats run from click-to-call, which lifts
click-through by 6--8\%, through interstitials, click-to-map, canvas and
expandable units, mobile search advertisements, in-app banners and rewarded video
--- and all of them convert better into a dynamic destination than a static one.
\end{summarybox}

%%%%%%%%%%%%%%%%%%%%%%%%%%%%%%%%%%%%%%%%%%%%%%%%%%%%%%%%%%%%%%%%%%%%%%%%%%%%%%
\chapter{Online Applications}
%%%%%%%%%%%%%%%%%%%%%%%%%%%%%%%%%%%%%%%%%%%%%%%%%%%%%%%%%%%%%%%%%%%%%%%%%%%%%%

\syllabusstrip{Online applications.}

\begin{learningoutcomes}
\begin{itemize}
\item The eight reasons apps are effective marketing tools.
\item The four app business models.
\item The three rules of app creation strategy.
\item The eleven features an app should incorporate.
\item Designing a mobile website, and why responsive design\ix{responsive design} superseded it.
\item Mobile analytics --- its significance and the seven metric groups.
\end{itemize}
\end{learningoutcomes}

\section{Why Apps Matter for Marketing}

\begin{keybox}[High Yield]
``How do online applications serve as effective tools for mobile marketing? What
are the key features that businesses should incorporate into their apps to ensure
they are maximizing marketing potential?''
\pyq{Improvement CIE, 28 Aug 2024 --- 10 M}
\end{keybox}

\begin{itemize}
\item \sfb{Permanent presence} --- the app icon sits on the home screen, creating
      continuous brand visibility.
\item \sfb{Push notification access} --- a direct, free, owned communication
      channel that does not depend on an algorithm.
\item \sfb{Rich behavioural data} --- in-app analytics reveal exactly how each
      user behaves, enabling deep personalisation.
\item \sfb{Higher engagement and conversion} --- app users typically browse more,
      convert more and spend more than mobile-web users.
\item \sfb{Loyalty and habit formation} --- points, streaks and saved preferences
      increase switching cost.
\item \sfb{Device capability access} --- camera, GPS, biometrics, wallet and
      contacts enable experiences the web cannot match.
\item \sfb{Offline availability} --- content and coupons remain usable without a
      connection.
\item \sfb{Frictionless repeat purchase} --- saved payment details and one-tap
      reorder.
\end{itemize}

\section{App Business Models}

\begin{itemize}
\item Free apps supported by advertising.
\item Paid apps supported by download fees.
\item Premium apps supported by in-app commerce.
\item Free apps supported by brands interested in connecting with customers.
\end{itemize}

\section{App Creation Strategy --- Three Rules}

\begin{frameworkbox}[The Three Rules]
\textbf{1. Make sure your app solves a problem.} The most effective apps solve
some problem for the user --- they facilitate a purchase, provide content, create
brand preference, or some combination. Analyse which of these your app will solve
and begin the design process from there.

\textbf{2. Get inside the mind of your user.} Understand how the user will engage
with the app --- will they be at home or in the office? Do they want information,
or simply to engage with the brand?

\textbf{3. Design with the end in mind.} Is the purpose to create brand
preference? To facilitate a financial transaction? To reduce customer churn\ix{churn}? Or
all of these?
\end{frameworkbox}

\section{Key Features an App Should Incorporate}

\begin{itemize}
\item \sfb{Fast onboarding} --- value visible before any sign-up demand.
\item \sfb{Personalised home screen} driven by behaviour.
\item \sfb{Search and filtering} that works.
\item \sfb{Push notification preferences} the user controls.
\item \sfb{In-app offers, loyalty and wallet.}
\item \sfb{One-tap checkout} with saved payment methods.
\item \sfb{Order tracking and history.}
\item \sfb{In-app support and chat.}
\item \sfb{Social sharing and referral.}
\item \sfb{Offline mode} for core content.
\item \sfb{Analytics instrumentation} on every meaningful event.
\end{itemize}

\section{Creating a Mobile Website}

With a smartphone you can browse any website available on a PC, but those sites
are generally unsuitable for a phone because they were not designed for one. The
screen is much smaller and behaviour differs --- a user may browse and read for
hours on a desktop, but on a phone will typically read only a small amount of
information over short periods. How the information is displayed, and how much of
it there is, must therefore be reconsidered.

\noindent When designing a mobile website, consider:

\begin{itemize}
\item You can only view one screen at a time --- design the navigation
      accordingly.
\item There is not much room for text, so don't use much.
\item Use large buttons for key calls to action.
\item Think about font usage --- make sure the important content really stands
      out.
\end{itemize}

\begin{keybox}[Why Responsive Design Superseded the Separate Mobile Site]
The modern equivalent of this advice is \keytermas{responsive
design}{responsive design} --- one codebase that adapts fluidly to every screen.
It is now preferred over maintaining a separate mobile site, because it avoids
duplicate content, splits no ranking authority and halves maintenance.
\end{keybox}

\section{Mobile Analytics}

\begin{keybox}[High Yield]
``Explain the significance of mobile analytics in mobile marketing campaigns. What
key metrics should marketers track to measure the success of their mobile
strategies?''
\pyq{SEE, Oct/Nov 2023, Q8b --- 8 M}
\end{keybox}

\begin{exambox}[Significance]
Mobile analytics is the only way to see behaviour \emph{inside} an app, where web
analytics cannot reach. It reveals which acquisition sources deliver users who
actually retain rather than merely install; it exposes where users abandon the
onboarding or checkout flow; it enables cohort analysis so that retention can be
compared across releases; it supports A/B testing\ix{A/B testing} of in-app experiences; and it
links marketing spend to lifetime value rather than to installs.
\end{exambox}

\begin{figure}[htbp]
\centering
\begin{tikzpicture}[font=\sffamily\footnotesize,
    grp/.style={draw=ocre,line width=0.85pt,fill=ocrepale,rounded corners=3pt,
                minimum height=1.25cm,text width=1.72cm,align=center,
                font=\sffamily\bfseries\scriptsize,text=slate,inner sep=3pt}]
  \node[grp] (g1) at (0,0)     {ACQUISITION};
  \node[grp] (g2) at (2.0,0)   {ACTIVATION};
  \node[grp] (g3) at (4.0,0)   {ENGAGEMENT};
  \node[grp] (g4) at (6.0,0)   {RETENTION};
  \node[grp,draw=greenok,fill=greenpale] (g5) at (8.0,0) {MONETISATION};
  \foreach \a/\b in {g1/g2,g2/g3,g3/g4,g4/g5}{
    \draw[-{Latex[length=1.9mm]},ocre,line width=0.85pt] (\a)--(\b);}

  \node[grp,draw=deepblue,fill=deepbluepale] (g6) at (10.15,0.75)
    {NOTIFICATION};
  \node[grp,draw=deepblue,fill=deepbluepale] (g7) at (10.15,-0.75)
    {TECHNICAL};
  \node[anchor=north,text=slatelight,font=\sffamily\scriptsize\itshape,
        align=center,text width=11cm] at (5.1,-1.0)
    {the five funnel groups run left to right; notification and technical health
     cut across all of them};
\end{tikzpicture}
\caption[The mobile analytics metric groups]%
        {Five funnel groups plus two cross-cutting ones. Judging a campaign on
         installs alone stops at the first box.}
\label{fig:mobanalytics}
\end{figure}

\begin{mmlongtable}{The seven mobile analytics metric groups}{tab:mobmetrics}%
{@{}P{0.22\textwidth}P{0.71\textwidth}@{}}
\rowcolor{ocre}\thd{Metric group} & \thd{Metrics}\\
\midrule
\endfirsthead
\rowcolor{ocre}\thd{Metric group} & \thd{Metrics}\\
\midrule
\endhead
\rlab{Acquisition} & Installs, cost per install, install source and attribution,
click-to-install rate\\
\rowcolor{tablealt}
\rlab{Activation} & Onboarding completion rate, time to first key action,
registration rate\\
\rlab{Engagement} & Daily and monthly active users, stickiness (daily divided by
monthly actives), session length, sessions per user, screens per session\\
\rowcolor{tablealt}
\rlab{Retention} & Day-1, Day-7 and Day-30 retention, churn rate, cohort
retention curves\\
\rlab{Monetisation} & Conversion rate, average revenue per user, average order
value, lifetime value, the ratio of lifetime value to customer acquisition cost\\
\rowcolor{tablealt}
\rlab{Notification performance} & Push opt-in rate, delivery rate, open rate,
opt-out rate\\
\rlab{Technical} & Crash rate, app load time, API latency, app store rating\\
\end{mmlongtable}

\begin{summarybox}
Apps are effective marketing tools because they give permanent home-screen
presence, owned push access, rich behavioural data, higher engagement and
conversion, loyalty and habit formation, device capability access, offline
availability and frictionless repeat purchase. Four business models fund them.
Three rules govern their creation: solve a real problem, understand how the user
will engage, and design with the end purpose in mind. Eleven features maximise
marketing potential, from fast onboarding to full analytics instrumentation. A
mobile site must be designed for one screen at a time, little text and large
buttons --- and responsive design has superseded the separate mobile site because
it avoids duplicate content and split authority. Mobile analytics is the only view
inside the app, and it is measured in seven groups: acquisition, activation,
engagement, retention, monetisation, notification performance and technical
health.
\end{summarybox}

%%%%%%%%%%%%%%%%%%%%%%%%%%%%%%%%%%%%%%%%%%%%%%%%%%%%%%%%%%%%%%%%%%%%%%%%%%%%%%
%  Unit IV appendices
%%%%%%%%%%%%%%%%%%%%%%%%%%%%%%%%%%%%%%%%%%%%%%%%%%%%%%%%%%%%%%%%%%%%%%%%%%%%%%

\startappendices

\chapter{Syllabus-to-Chapter Concordance}

This appendix exists so that coverage can be \emph{audited} rather than assumed.
The left-hand column reproduces the Unit~IV syllabus descriptor phrase by phrase,
in the order the official document prints it. The right-hand column gives the
chapter and section that carries it. Every phrase resolves to a location, and no
numbered chapter of this volume contains material that is not traceable to a
phrase in this table.

\begin{mmlongtable}{Concordance for Unit IV}{tab:concordance4}%
{@{}P{0.36\textwidth}P{0.57\textwidth}@{}}
\rowcolor{ocre}\thd{Syllabus phrase} & \thd{Where it is covered}\\
\midrule
\endfirsthead
\rowcolor{ocre}\thd{Syllabus phrase} & \thd{Where it is covered}\\
\midrule
\endhead
\multicolumn{2}{@{}l}{\bfseries\color{ocredark} First half --- Web Design}\\
\midrule
\rlab{Web design, optimization of websites} & Chapter~1 entire. What web design is
\S\,1.1; UI and UX \S\,1.2; website structure \S\,1.3; website optimization
defined \S\,1.4; the eight techniques \S\,1.5; image optimization \S\,1.6;
testing and optimisation \S\,1.7.\\
\rowcolor{tablealt}
\rlab{Publishing a basic website} & Chapter~2 entire. The nine official steps
\S\,2.1; each step expanded, including the pre-launch technical setup \S\,2.2.\\
\rlab{User-centered Design and Website Optimization} & Chapter~3 entire. What UCD
is \S\,3.1; the seven principles \S\,3.2; the UCD process \S\,3.3; what visitors
look for \S\,3.4. Website optimization itself is carried in \S\,1.4--1.7.\\
\rowcolor{tablealt}
\rlab{Design Principles and Website Copy} & Chapter~4 entire. The core design
principles \S\,4.1; Krug's first law and the usability attributes \S\,4.2; how
people use the web \S\,4.3; billboard design \S\,4.4; copy principles \S\,4.5;
design and copy as one system \S\,4.6.\\
\rlab{Website Metrics \& Developing Insight} & Chapter~5 entire. Key metrics
\S\,5.1; how metrics aid optimisation \S\,5.2; from data to decision \S\,5.3;
user behaviour analysis \S\,5.4.\\
\midrule
\multicolumn{2}{@{}l}{\bfseries\color{ocredark} Second half --- Mobile Marketing}\\
\midrule
\rowcolor{tablealt}
\rlab{Difference between mobile advertising and marketing} & Chapter~6 entire.
What mobile marketing is \S\,6.1; advertising against marketing \S\,6.2; mobile
web against app \S\,6.3; benefits \S\,6.4; limitations \S\,6.5; the seven
campaign characteristics \S\,6.6; geo-targeting \S\,6.7.\\
\rlab{Utilizing mobile marketing for sales promotions} & Chapter~7 entire. Push
notifications \S\,7.1; ten promotion techniques \S\,7.2; advertising formats
\S\,7.3.\\
\rowcolor{tablealt}
\rlab{Online applications} & Chapter~8 entire. Why apps matter \S\,8.1; business
models \S\,8.2; the three creation rules \S\,8.3; the feature checklist \S\,8.4;
creating a mobile website \S\,8.5; mobile analytics \S\,8.6.\\
\end{mmlongtable}

\begin{keybox}[Two Notes on Coverage]
\textbf{Deliberate overlap with Volume~1.} Krug's first law, the seven usability\ix{usability}
attributes and the landing-page material appear in Unit~I under ``Optimize
Website UX \& Landing Pages'' and again here under ``Design Principles''. Both
descriptors name them, so both volumes carry them. Chapter~4 gives the design-side
treatment; Volume~1 Chapter~9 gives the content-marketing side.

\smallskip
\textbf{No off-syllabus appendix.} As in Volume~3, every topic in the source
material for Unit~IV maps to a descriptor phrase, so nothing had to be set aside.
\end{keybox}

\chapter{Previous-Year Question Bank --- Unit IV}

Every question in this appendix is reproduced from an actual RVCE Marketing
Management Continuous Internal Evaluation or Semester End Examination paper for
this course. Answers and pointers follow the official schemes of valuation where
those were released.

\section*{Part A --- Two-Mark Questions}
\addcontentsline{toc}{section}{Part A --- Two-Mark Questions}

\begin{enumerate}[leftmargin=2.2em,itemsep=1.1ex]

\item \textbf{What is Website Optimization?} \pyqr{CIE-III, Jun 2026 --- 2 M}\\
      The process of improving website performance, speed, usability and search
      visibility in order to increase traffic and conversions.
      \hfill\emph{See} \S\,1.4.

\item \textbf{Mention any two principles of website design.}
      \pyqr{CIE-III, Jun 2026 --- 2 M}\\
      Simplicity and consistency --- these principles improve readability and user
      experience. \hfill\emph{See} \S\,4.1.

\item \textbf{Differentiate Mobile Advertising and Mobile Marketing. / What is the
      key difference between mobile advertising and mobile marketing? / Define
      mobile advertising in one sentence. How does it differ from mobile
      marketing?}
      \pyqr{CIE-III, Jun 2026; Improvement CIE, 4 Jun 2025; Improvement CIE, 28 Aug 2024 --- 2 M}\\
      Mobile advertising focuses mainly on paid promotional messages delivered to
      mobile devices. Mobile marketing includes overall customer engagement
      through SMS, apps, social media and mobile websites. In one line: mobile
      advertising is a paid subset of the broader mobile marketing strategy.
      \hfill\emph{See} \S\,6.2.

\item \textbf{Mention two key principles of User-Centered Design in website
      development.} \pyqr{Improvement CIE, 4 Jun 2025 --- 2 M}\\
      (1) Design based on an explicit understanding of the users, their tasks and
      their environment. (2) Iterative design that is continuously refined
      through user involvement and user-centred evaluation.
      \hfill\emph{See} \S\,3.1.

\item \textbf{How can website metrics help in optimizing a website for marketing
      effectiveness?} \pyqr{Improvement CIE, 4 Jun 2025 --- 2 M}\\
      They locate problems precisely, rank fixes by impact, prove or disprove
      changes through testing, reallocate budget to converting channels, and
      expose segment-specific failures hidden by averages.
      \hfill\emph{See} \S\,5.2.

\item \textbf{State two key elements that visitors commonly expect to see on a
      well-optimized landing page.} \pyqr{Improvement CIE, 4 Jun 2025 --- 2 M}\\
      A clear benefit-led headline matching the advertisement or link that brought
      them there, and a single prominent call-to-action button. Also acceptable:
      trust signals such as testimonials or ratings. \hfill\emph{See} \S\,3.4.

\item \textbf{How does image optimization contribute to overall website
      performance?} \pyqr{Improvement CIE, 28 Aug 2024 --- 2 M}\\
      It reduces page weight and load time, improving Core Web Vitals\ix{Core Web Vitals} and
      lowering bounce\ix{bounce} rate; improves the mobile experience on slow connections;
      improves SEO through descriptive file names and ALT text; and improves
      accessibility. \hfill\emph{See} \S\,1.6.

\item \textbf{List any four key metrics used to measure the success of a website.}
      \pyqr{Improvement CIE, 28 Aug 2024 --- 2 M}\\
      Traffic (users or sessions), bounce rate, average session duration and
      conversion rate. Also acceptable: pages per session, page load time,
      traffic sources, new against returning visitors.
      \hfill\emph{See} \S\,5.1.

\item \textbf{How can push notifications be used in mobile marketing to enhance
      sales promotions?} \pyqr{Improvement CIE, 28 Aug 2024 --- 2 M}\\
      They deliver time-critical offers instantly to the installed base, create
      urgency for flash sales, recover abandoned carts, trigger location-based
      coupons when a user nears a store, and re-engage lapsed users --- all
      through an owned channel that no algorithm throttles.
      \hfill\emph{See} \S\,7.1.

\item \textbf{How can user-friendly design impact the overall effectiveness of a
      website?} \pyqr{Improvement CIE, 28 Aug 2024 --- 2 M}\\
      It reduces cognitive effort so that visitors can complete their task, which
      lowers bounce rate, increases time on site and pages per session, raises
      conversion rate, builds trust in the brand, encourages return visits, and
      --- because search engines now mimic human preferences --- indirectly
      improves organic ranking. \hfill\emph{See} \S\,4.2.

\item \textbf{Why is it important to integrate SEO in website development?}
      \pyqr{SEE, Jun/Jul 2025 --- 2 M}\\
      Architecture, URLs, rendering method, speed, heading structure and analytics
      are all fixed at build time. Retro-fitting SEO is expensive, forces URL
      changes that risk losing ranking authority, and often leaves staging
      \texttt{noindex} directives live in production.
      \hfill\emph{See} \S\,2.2, and Volume~2 \S\,8.2.

\item \textbf{What is the role of geo-targeting in mobile marketing?}
      \pyqr{SEE, Jun/Jul 2025 --- 2 M}\\
      Geo-targeting delivers content and offers based on the user's location,
      making messages locally relevant and timely --- driving footfall through
      geo-fenced offers, enabling ``near me'' discovery and click-to-map
      directions, supporting location-specific pricing and inventory, and
      improving ROI by eliminating spend outside the serviceable area.
      \hfill\emph{See} \S\,6.7.

\item \textbf{List any two aspects of branding and consistency in web design.}
      \pyqr{SEE, Oct/Nov 2023 --- 2 M}\\
      (1) A consistent visual identity --- the same logo, colour palette,
      typography, imagery style and button design applied on every page. (2) A
      consistent tone of voice and terminology in all copy, including microcopy,
      error messages and CTA labels, so that the brand reads as one coherent
      personality. \hfill\emph{See} \S\,4.1.

\item \textbf{List any two differences between Mobile Web marketing and Mobile App
      marketing.} \pyqr{SEE, Oct/Nov 2023 --- 2 M}\\
      (1) Mobile web is accessed through a browser with no installation and
      therefore has broader reach, whereas an app requires download and
      installation and reaches a narrower but far more engaged audience. (2) Apps
      have full push notification and device-feature access --- camera, GPS,
      biometrics --- and can work offline, whereas mobile web has limited push
      capability and device access. \hfill\emph{See} \S\,6.3.

\item \textbf{What is the impact of user behavior analysis on website
      optimization?} \pyqr{SEE, Oct/Nov 2023 --- 2 M}\\
      It shows what users actually do rather than what is assumed --- identifying
      friction points, ignored content, invisible calls to action, rage clicks and
      true scroll depth\ix{scroll depth} --- and generates the hypotheses that A/B testing then
      validates. Without it, optimisation is guesswork.
      \hfill\emph{See} \S\,5.4.

\end{enumerate}

\section*{Part B --- Eight- and Ten-Mark Questions}
\addcontentsline{toc}{section}{Part B --- Eight- and Ten-Mark Questions}

\begin{enumerate}[leftmargin=2.2em,itemsep=1.4ex]

\item \textbf{Explain the concept of Web Design and the importance of website
      optimization.} \pyqr{CIE-III, Jun 2026 --- 10 M}\\
      Define web design covering layout, colour, typography, graphics, navigation
      and responsiveness. Then give the six importance points, the four types of
      optimization --- technical, content, SEO, mobile --- and the three benefits.
      \hfill\emph{See} \S\,1.1 and \S\,1.4.

\item \textbf{A start-up company notices that visitors leave their website within
      a few seconds. As a web designer, explain how you would optimize the website
      to improve user engagement and conversion rate.}
      \pyqr{CIE-III, Jun 2026 --- 10 M}\\
      Give the eight numbered techniques --- speed, responsive design\ix{responsive design}, navigation,
      content quality, strong calls to action, SEO, bounce-rate reduction and
      metric monitoring --- each with its concrete sub-actions. Strengthen the
      answer by opening with a diagnosis step (heatmaps and session recordings to
      find \emph{why} they leave) and closing with A/B testing\ix{A/B testing} to validate the
      fixes. \hfill\emph{See} \S\,1.5.

\item \textbf{A college plans to publish its departmental website for admissions
      and student activities. Explain the complete process involved in designing
      and publishing the website.} \pyqr{CIE-III, Jun 2026 --- 10 M}\\
      Give the nine official steps --- requirement analysis, website design,
      content development, domain registration, web hosting, website development,
      testing, publishing, maintenance --- expanded with the practitioner detail,
      including the pre-launch technical setup. \hfill\emph{See} Chapter~2.

\item \textbf{Discuss what website visitors typically look for when they land on a
      business webpage. Suggest effective design and content strategies to meet
      those expectations.} \pyqr{Improvement CIE, 4 Jun 2025 --- 10 M}\\
      Use the ten visitor questions with the paired delivery strategy for each,
      then add the design principles and the copy principles as the systematic
      response. \hfill\emph{See} \S\,3.4, \S\,4.1 and \S\,4.5.

\item \textbf{Design a user-centered website layout for a new e-commerce startup
      targeting young urban professionals. How would you ensure optimal user
      experience and engagement? Include key design principles, optimization
      strategies, and website copy suggestions.}
      \pyqr{Improvement CIE, 4 Jun 2025 --- 10 M}\\
      \emph{Research first:} persona --- 25 to 35, urban, mobile-first, time-poor,
      price-aware but convenience-driven; primary task --- find and buy quickly on
      a phone during a commute.
      \emph{Layout:} sticky header with search and cart; a hero section with one
      benefit headline and one call to action\ix{call to action}; a three-column category grid; a
      social-proof band with ratings; a personalised ``recommended for you'' row;
      a simple footer with trust and policy links.
      \emph{Key design principles applied:} mobile-first responsive layout; visual
      hierarchy placing the search bar and primary CTA highest; simplicity with
      generous whitespace; consistency in button style and colour; one accent
      colour reserved for actions; 16-pixel or larger type; tap targets of at
      least 44 pixels; accessibility through ALT text and contrast.
      \emph{Optimization strategies:} compressed modern-format images and lazy
      loading for a sub-three-second load; one-page guest checkout with a
      four-field form and saved payment; a progress indicator in checkout; an
      exit-intent offer; heatmaps and session recordings to locate friction; A/B
      tests on the hero headline and CTA.
      \emph{Copy suggestions:} headline ``Everything for your week, delivered by
      tomorrow''; sub-line ``Free delivery over \rupee 499. Returns in one tap.'';
      CTA ``Shop today's deals'' rather than ``Submit''; microcopy under the form
      reading ``Takes 30 seconds --- no account needed''; benefit-led product copy
      with specific claims.
      \emph{Measurement:} conversion rate, checkout completion, mobile bounce rate
      and average order value. \hfill\emph{See} Chapters~3, 4 and~5.

\item \textbf{A travel agency has created a landing page for a limited-time offer
      on holiday packages. Despite a high number of clicks on the ads leading to
      the landing page, the conversion rate is disappointingly low. Evaluate the
      possible reasons for the low conversion rate. Suggest optimizations for key
      elements such as headlines, call-to-action buttons, and visuals.}
      \pyqr{Improvement CIE, 28 Aug 2024 --- 10 M}\\
      \emph{Possible reasons:} poor message match\ix{message match} --- the advertisement promised
      something the page does not repeat; a vague headline that does not restate
      the offer; the offer and price are below the fold; too many competing calls
      to action, or none that is visually dominant; a long enquiry form demanding
      unnecessary details; no urgency despite the offer being time-limited;
      missing trust signals --- no reviews, no certifications, no cancellation
      policy; slow page load, especially on mobile; generic stock imagery that does
      not show the actual destinations; hidden total cost or unclear inclusions;
      and navigation links that let visitors wander away from the page.

      \emph{Optimisations --- headlines:} restate the exact advertisement promise
      with specifics, such as ``Kerala backwaters, 4 nights from \rupee 18,999 ---
      book by Sunday'', and add a sub-headline listing what is included.
      \emph{CTA buttons:} one dominant, high-contrast accent-coloured button
      repeated above and below the fold, with active benefit copy such as
      ``Reserve my package'' instead of ``Submit'', plus reassurance microcopy
      (``Free cancellation up to 7 days'').
      \emph{Visuals:} replace stock photography with real destination and customer
      photographs, add a short itinerary video, and show a countdown timer for the
      deadline.
      \emph{Additional:} cut the form to name, phone and travel date; add
      testimonials and ratings adjacent to the CTA; remove the site navigation from
      the landing page; compress images for a sub-three-second mobile load; and
      A/B test the headline and CTA.
      \emph{Measurement:} conversion rate, form-start to form-completion ratio,
      scroll depth, and session recordings to verify the fix.
      \hfill\emph{See} \S\,3.4, \S\,4.5 and \S\,5.3.

\item \textbf{How do online applications serve as effective tools for mobile
      marketing? What are the key features that businesses should incorporate into
      their apps to ensure they are maximizing marketing potential?}
      \pyqr{Improvement CIE, 28 Aug 2024 --- 10 M}\\
      The eight reasons apps are effective, the four business models, the three
      app-creation rules, and the eleven-feature checklist.
      \hfill\emph{See} Chapter~8.

\item \textbf{A luxury cosmetics brand is planning to launch a flash sale
      exclusively through mobile devices. How would you design a mobile marketing
      campaign that not only drives awareness of the sale but also maximizes
      conversion rates during the limited time?}
      \pyqr{Improvement CIE, 28 Aug 2024 --- 10 M}\\
      \emph{Pre-launch (awareness, seven days out):} teaser push notifications and
      SMS to the segmented installed base; Reels and short-video influencer
      previews; an in-app countdown banner; a ``notify me'' opt-in that captures
      intent and grants push permission; paid social and mobile display
      retargeting to past purchasers and lookalikes.
      \emph{Launch (conversion window):} a timed push notification at each
      segment's peak activity hour, deep-linked straight to the sale screen;
      app-exclusive early access one hour before the website, to drive installs; a
      visible countdown timer on the sale screen; personalised product rows based
      on browsing and purchase history; limited-stock indicators to create
      scarcity; one-tap checkout with saved cards and wallet support.
      \emph{During (maximising conversion):} an abandoned-cart push within 20
      minutes; geo-fenced offers near flagship stores with click-to-map; live chat
      support in-app; free-shipping threshold prompts to raise average order value.
      \emph{Post-launch:} thank-you and review-request push; a win-back offer to
      browsers who did not buy; and a wallet pass for the next drop.
      \emph{Technical requirements:} sub-three-second load, server capacity for the
      traffic spike, and a fully responsive mobile-first design.
      \emph{Measurement:} push open rate, sale-screen sessions, add-to-cart rate,
      checkout completion, conversion rate, average order value, revenue per push,
      opt-out rate, and app installs during the window.
      \hfill\emph{See} Chapters~7 and~8.

\item \textbf{Assume you are launching a promotional campaign for a new product
      using mobile marketing. Prepare a detailed plan including mobile advertising
      channels, push notifications, app-based promotions, and integration with
      social media.} \pyqr{Improvement CIE, 4 Jun 2025 --- 10 M}\\
      \emph{Objectives and audience:} define a SMART\ix{SMART} objective and the mobile
      persona, including device, active hours and preferred platforms.
      \emph{Mobile advertising channels:} mobile search advertisements on
      high-intent queries with location extensions; in-app banner and native
      advertisements on relevant apps; rewarded video for reach; interstitials at
      natural app breakpoints; click-to-call advertisements for enquiry-driven
      products; click-to-map advertisements to drive footfall; and expandable
      rich-media advertisements for the brand story.
      \emph{Push notifications:} a sequenced plan --- teaser, launch, reminder,
      last chance --- each segmented, deep-linked, frequency-capped and A/B tested,
      with permission requested only after the user has experienced value.
      \emph{App-based promotions:} app-exclusive early access and pricing, in-app
      coupons, gamified spin-to-win, loyalty points, referral codes, wallet passes,
      and one-tap checkout.
      \emph{Social media integration:} Reels and Stories with the campaign
      hashtag; influencer seeding; shoppable posts and product tags; a messaging
      broadcast for order updates and offers; user-generated content\ix{user-generated content} contests; and
      retargeting audiences built from app and site events.
      \emph{Cross-channel integration:} QR codes on packaging and print linking to
      the mobile landing page; consistent creative and offer across all
      touchpoints; ``think mobile first'', with mobile as the connective tissue
      between all media.
      \emph{Measurement:} installs and cost per install, push open rate, in-app
      conversion rate, coupon redemption, return on advertising spend, and lifetime
      value by acquisition source. \hfill\emph{See} Chapters~6, 7 and~8.

\item \textbf{Discuss the relationship between design principles and website copy.
      How can visual design elements and well-crafted content work together to
      enhance user experience and achieve website objectives? Include examples of
      effective design and copy integration.}
      \pyqr{Improvement CIE, 28 Aug 2024 --- 10 M}\\
      The principle-by-principle table, followed by the worked examples of
      successful integration --- hero section, pricing table, testimonial block ---
      and of failure: a beautiful layout with vague copy, or strong copy set in
      unreadable type. \hfill\emph{See} \S\,4.6.

\item \textbf{What is the role of user Interface and user Experience designs in web
      design for marketing? How can businesses ensure a positive user Journey on
      their websites?} \pyqr{SEE, Oct/Nov 2023, Q7a --- 8 M}\\
      The UI and UX comparison table and the eight ways to ensure a positive user
      journey. \hfill\emph{See} \S\,1.2.

\item \textbf{Describe the importance of testing and optimization in web design for
      Marketing. What tools and methods can be used to analyze and improve Website
      performance?} \pyqr{SEE, Oct/Nov 2023, Q7b --- 8 M}\\
      The case for testing and the method-and-tool table.
      \hfill\emph{See} \S\,1.7.

\item \textbf{Discuss briefly on Push Notification in Mobile Marketing. Explain
      how can businesses use them to engage and retain Mobile App users? Provide
      examples of effective push notification strategies.}
      \pyqr{SEE, Oct/Nov 2023, Q8a --- 8 M}\\
      The definition, the seven engagement and retention uses, and the eight
      effective strategies --- with examples such as a 20-minute abandoned-cart
      reminder, a geo-fenced coupon on store approach, and a 30-day win-back
      message. \hfill\emph{See} \S\,7.1.

\item \textbf{Explain the significance of mobile analytics in mobile marketing
      campaigns. What key metrics should marketers track to measure the success of
      their mobile strategies?} \pyqr{SEE, Oct/Nov 2023, Q8b --- 8 M}\\
      The significance argument and the seven metric groups, from acquisition
      through to technical health. \hfill\emph{See} \S\,8.6.

\item \textbf{Evaluate the effectiveness of mobile marketing in facing market
      changes due to rapid technological advancements.}
      \pyqr{SEE, Jun/Jul 2025, Q8a --- 8 M}\\
      \emph{Where mobile marketing is effective:} it is inherently adaptive
      because it is digital and measurable --- campaigns can be tested and
      adjusted continuously; it reaches consumers wherever they are, at any hour;
      it spans the entire funnel from discovery to purchase; it costs less than
      traditional advertising; it produces first-party behavioural data that
      survives the loss of third-party cookies; and it enables geo-targeting,
      personalisation and instant response that no traditional channel can match.

      \emph{Where it is strained by technological change:} platform fragmentation
      across devices, screen sizes and operating systems makes a single campaign
      difficult; privacy shifts --- app tracking transparency, cookie deprecation,
      GDPR and do-not-disturb regulation --- have degraded attribution and
      targeting; ad fatigue and rising uninstall rates limit push frequency; and
      the pace of format change forces continual re-skilling.

      \emph{Evaluation:} mobile marketing remains the most resilient channel
      precisely because its measurability lets it adapt faster than the disruption
      --- but only for organisations that build an owned first-party asset, namely
      an app and a permission-based list, rather than renting reach from
      platforms. \hfill\emph{See} \S\,6.4, \S\,6.5 and \S\,6.6.

\item \textbf{Create a strategy to use new digital channels like mobile apps,
      chatbots, and social media integration to strengthen a brand's global online
      presence.} \pyqr{SEE, Jun/Jul 2025, Q8b --- 8 M}\\
      \emph{Mobile app as the owned hub:} a single global app with localised
      content, currency and language; personalised home screens; push notification
      consent captured after first value; loyalty and wallet integration; offline
      access; and full event instrumentation feeding a global analytics view.

      \emph{Chatbots as the always-on layer:} deployed on the website, in the app
      and on messaging platforms; handling FAQs, order tracking, returns and
      product discovery in multiple languages across all time zones; escalating to
      human agents with full context; and capturing conversational data that feeds
      product and content decisions.

      \emph{Social media integration:} a hub-and-spoke model with a global brand
      account and regional handles; platform-native content rather than
      cross-posting; shoppable posts and product tags linking back to the app;
      regional influencer partnerships for cultural credibility; and social
      listening for share of voice\ix{share of voice} and sentiment by market.

      \emph{Cross-channel integration:} one customer profile unified across app,
      chat, social, e-mail and web, so that the conversation continues seamlessly
      --- the omnichannel synchronisation principle of Volume~2 \S\,1.3.

      \emph{Governance:} a central brand system with regional flexibility;
      consistent visual identity and tone; and compliance with local privacy and
      messaging regulation.

      \emph{Measurement:} a global dashboard tracking reach, engagement, app
      installs and retention, chatbot resolution rate, cross-channel conversion and
      lifetime value by market. \hfill\emph{See} Chapters~6 and~8.

\end{enumerate}

\section*{How Unit IV Is Examined}
\addcontentsline{toc}{section}{How Unit IV Is Examined}

\begin{keybox}[The Pattern Across Papers]
In the Semester End Examination, Unit~IV is Questions~7 and~8 with internal
choice --- 16 marks in two eight-mark sub-divisions --- plus two or three Part~A
objective questions. The third Continuous Internal Evaluation test and the
Improvement CIE both draw heavily on it.

\smallskip
The highest-probability short question by a wide margin is the \textbf{mobile
advertising against mobile marketing} distinction, which has appeared in three
separate papers. The highest-probability ten-mark topics are the
\textbf{website optimisation scenario} (visitors leaving within seconds), the
\textbf{nine-step publishing process}, \textbf{what visitors look for} on a
business page, the \textbf{design-and-copy relationship}, \textbf{push
notifications}, and \textbf{online applications}.

\smallskip
Note that Unit~IV questions are heavily scenario-based --- a start-up, a college
department, a travel agency, a cosmetics brand, an e-commerce startup. Answer them
by naming the framework and then applying it to the specific business, not by
discussing web design in general.
\end{keybox}

\chapter{Glossary of Key Terms}

\begin{mmlongtable}{Key terms for rapid revision}{tab:glossary4}%
{@{}P{0.26\textwidth}P{0.67\textwidth}@{}}
\rowcolor{ocre}\thd{Term} & \thd{One-line meaning}\\
\midrule
\endfirsthead
\rowcolor{ocre}\thd{Term} & \thd{One-line meaning}\\
\midrule
\endhead
\rlab{Above the fold} & The area of a page visible without scrolling\\
\rowcolor{tablealt}
\rlab{ASO} & App store optimisation --- ranking an app in app store search\\
\rlab{Click-to-Call} & A mobile advertisement format with a tappable phone number;
lifts click-through by 6--8\%\\
\rowcolor{tablealt}
\rlab{Click-to-Map} & A geo-targeted format that opens a map to the nearest
store\\
\rlab{Core Web Vitals} & The page-experience metrics LCP, INP and CLS\\
\rowcolor{tablealt}
\rlab{CPI / LTV} & Cost per install / customer lifetime value\\
\rlab{DAU / MAU} & Daily and monthly active users; their ratio is
``stickiness''\\
\rowcolor{tablealt}
\rlab{Domain registration} & Buying the site's address\\
\rlab{Geo-fencing} & Triggering an action when a user enters a defined geographic
boundary\\
\rowcolor{tablealt}
\rlab{Geo-targeting} & Delivering content or offers based on the user's geographic
location\\
\rlab{Happy talk} & Empty welcome text that says nothing; should be deleted\\
\rowcolor{tablealt}
\rlab{Interstitial} & A full-screen advertisement displayed within an app\\
\rlab{Krug's first law} & ``Don't make me think'' --- the interface must be
self-evident\\
\rowcolor{tablealt}
\rlab{M-commerce} & Commercial transactions conducted through a wireless
internet-enabled device\\
\rlab{Mobile advertising} & The paid promotional subset of mobile marketing\\
\rowcolor{tablealt}
\rlab{Mobile marketing} & A multi-channel strategy reaching users on mobile
devices\\
\rlab{Push notification} & A short message sent by an app to the device screen,
even when the app is closed\\
\rowcolor{tablealt}
\rlab{Responsive design} & One layout that adapts fluidly to mobile, tablet and
desktop\\
\rlab{Satisficing} & Choosing the first option that plausibly works rather than
the best one\\
\rowcolor{tablealt}
\rlab{UI / UX} & The visual interactive surface / the whole quality of the user's
journey\\
\rlab{User-centred design} & Design driven by explicit user understanding and
validated iteratively by user testing\\
\rowcolor{tablealt}
\rlab{Visual hierarchy} & Making important elements visually prominent in
proportion to their importance\\
\rlab{Web design} & Creating visually appealing, user-friendly websites --- layout,
colour, typography, graphics, navigation, responsiveness\\
\rowcolor{tablealt}
\rlab{Web hosting} & Renting the server that stores the site's files\\
\rlab{Website copy} & The written content of a page --- headlines, body, microcopy,
CTA text\\
\rowcolor{tablealt}
\rlab{Website optimization} & Improving website performance, speed, usability and
search visibility to increase traffic and conversions\\
\end{mmlongtable}

\chapter{Framework Quick-Reference Sheet}

Everything in this volume that is a numbered list, in one place, for the night
before the examination.

\begin{mmlongtable}{Framework quick reference}{tab:quickref4}%
{@{}P{0.28\textwidth}P{0.65\textwidth}@{}}
\rowcolor{ocre}\thd{Framework} & \thd{The sequence}\\
\midrule
\endfirsthead
\rowcolor{ocre}\thd{Framework} & \thd{The sequence}\\
\midrule
\endhead
\rlab{Web design defined} & Layout, colour, typography, graphics, navigation,
responsiveness\\
\rowcolor{tablealt}
\rlab{UI against UX} & Definition, scope, deliverables, marketing role, failure
symptom\\
\rlab{Positive user journey (eight actions)} & Map the journey, shallow
navigation, obvious primary action, fewest steps, feedback, error and empty
states, speed, test with real users\\
\rowcolor{tablealt}
\rlab{Basic pages} & Home, about, product and service, landing, blog, contact,
trust\\
\rlab{Website optimization --- importance} & Speed, experience, ranking, bounce,
conversion, mobile compatibility\\
\rowcolor{tablealt}
\rlab{Four types of optimization} & Technical, content, SEO, mobile\\
\rlab{Eight optimisation techniques} & Speed \ra{} responsive design \ra{}
navigation \ra{} content quality \ra{} strong CTAs \ra{} SEO \ra{} bounce
reduction \ra{} metric monitoring\\
\rowcolor{tablealt}
\rlab{Image optimization --- four gains} & Speed, mobile experience, SEO,
accessibility\\
\rlab{Testing methods} & A/B, usability, heatmaps, recordings, funnel analysis,
performance, cross-device, accessibility, crawl\\
\rowcolor{tablealt}
\rlab{Nine publishing steps} & Requirement analysis \ra{} design \ra{} content
\ra{} domain \ra{} hosting \ra{} development \ra{} testing \ra{} publishing \ra{}
maintenance\\
\rlab{UCD --- seven principles} & Explicit user understanding, continuous
involvement, evaluation-driven, iterative, whole experience, multidisciplinary,
accessible\\
\rowcolor{tablealt}
\rlab{UCD process} & Understand context \ra{} specify requirements \ra{} produce
solutions \ra{} evaluate \ra{} iterate\\
\rlab{Ten visitor questions} & Right place, what's in it for me, trust, cost, next
step, findability, speed, mobile, human contact, safety\\
\rowcolor{tablealt}
\rlab{Twelve design principles} & Simplicity, consistency, visual hierarchy,
balance and alignment, contrast, whitespace, typography, colour, responsiveness,
accessibility, navigation clarity, speed\\
\rlab{Seven usability attributes} & Useful, learnable, memorable, effective,
efficient, desirable, delightful\\
\rowcolor{tablealt}
\rlab{Billboard design rules} & Visual hierarchy, omit needless words, obvious
clickables, breadcrumbs and street signs\\
\rlab{Copy principles} & Benefit-led, scannable, front-loaded, customer's
language, specific, active, concise, one idea per section, explicit CTA, objections
answered in place\\
\rowcolor{tablealt}
\rlab{Design serving copy} & Hierarchy, typography, whitespace, contrast,
consistency, imagery, responsiveness\\
\rlab{Four metrics to name} & Traffic, bounce rate, conversion rate, average
session duration\\
\rowcolor{tablealt}
\rlab{Data to decision} & Observe \ra{} segment \ra{} diagnose \ra{} hypothesise
\ra{} test \ra{} decide \ra{} re-measure\\
\rlab{Mobile advertising against marketing} & Scope, nature, channels, objective,
duration, interaction, measurement --- advertising is the paid subset\\
\rowcolor{tablealt}
\rlab{Mobile web against app} & Access, reach, discoverability, engagement depth,
push, offline, device features, cost\\
\rlab{Five mobile limitations} & Platform diversity, privacy, small-screen
navigation, ad fatigue, regulation\\
\rowcolor{tablealt}
\rlab{Seven campaign characteristics} & Mobile first, multi-screen, full spectrum,
integrate with traditional, work across screens, whole funnel, test your way\\
\rlab{Push uses} & Immediacy, urgency, personalisation, cart recovery,
geo-triggered, re-engagement, loyalty\\
\rowcolor{tablealt}
\rlab{Eight push strategies} & Segment, time correctly, cap frequency, one benefit
with a verb, rich media, deep-link, ask after value, A/B test and watch opt-out\\
\rlab{Ten promotion techniques} & SMS, coupons, in-app exclusives, push flash
sales, location offers, QR codes, gamification, wallet passes, referrals, early
access\\
\rowcolor{tablealt}
\rlab{Mobile ad formats} & Click-to-call, interstitial, click-to-map, canvas and
expandable, mobile search, in-app banner and native, rewarded video\\
\rlab{Eight reasons apps work} & Permanent presence, push access, behavioural
data, higher engagement, loyalty, device capability, offline, frictionless
repurchase\\
\rowcolor{tablealt}
\rlab{Three app rules} & Solve a problem, get inside the user's mind, design with
the end in mind\\
\rlab{Seven mobile metric groups} & Acquisition, activation, engagement,
retention, monetisation, notification performance, technical\\
\end{mmlongtable}


%%%%%%%%%%%%%%%%%%%%%%%%%%%%%%%%%%%%%%%%%%%%%%%%%%%%%%%%%%%%%%%%%%%%%%%%%%%%%%
\backmatter
\printindex

\end{document}
